% ============================================================
% FINAL PROJECT REPORT --- ABAP19 Configuration Management
% FPT University Capstone Project --- Group GSP26
% Format: Times New Roman 12pt, 1.5 spacing, A4
% Margins: Top 2.5cm, Bottom 2.5cm, Left 3.5cm, Right 2cm
% ============================================================
\documentclass[12pt,a4paper]{report}

% --- Encoding ---
\usepackage[utf8]{inputenc}
\usepackage[T1]{fontenc}
\usepackage{mathptmx}              % Times New Roman (text + math)

% --- Page layout ---
\usepackage[top=2.5cm,bottom=2.5cm,left=3.5cm,right=2cm]{geometry}
\usepackage{setspace}
\onehalfspacing                     % 1.5 line spacing

% --- Graphics & Color ---
\usepackage{graphicx}
\usepackage[table]{xcolor}
\usepackage{float}

% --- Tables ---
\usepackage{longtable}
\usepackage{tabularx}
\usepackage{array}
\usepackage{multirow}
\usepackage{makecell}
\usepackage{colortbl}

% --- Lists ---
\usepackage{enumitem}

% --- Headers & footers ---
\usepackage{fancyhdr}
\pagestyle{fancy}
\fancyhf{}
\fancyfoot[C]{\thepage}
\renewcommand{\headrulewidth}{0pt}
\renewcommand{\footrulewidth}{0pt}

% --- Heading styles (FPT requirements) ---
\usepackage{titlesec}

% Chapter heading: not used directly; we use \section as top level
% But report class needs chapter for TOC structure.
% We make chapter print nothing visible (content starts with \section).
\titleformat{\chapter}[display]
  {\normalfont\fontsize{16}{20}\bfseries}
  {}{0pt}{}
\titlespacing*{\chapter}{0pt}{20pt}{15pt}

% Section = "1." --- 16pt bold
\titleformat{\section}[block]
  {\normalfont\fontsize{16}{20}\bfseries}
  {\thesection.}{0.5em}{}
\titlespacing*{\section}{0pt}{18pt}{8pt}

% Subsection = "1.1." --- 14pt bold
\titleformat{\subsection}[block]
  {\normalfont\fontsize{14}{17}\bfseries}
  {\thesubsection.}{0.5em}{}
\titlespacing*{\subsection}{0pt}{14pt}{6pt}

% Subsubsection = "1.1.1." --- 12pt bold (NO italic)
\titleformat{\subsubsection}[block]
  {\normalfont\fontsize{12}{15}\bfseries}
  {\thesubsubsection.}{0.5em}{}
\titlespacing*{\subsubsection}{0pt}{10pt}{4pt}

% Paragraph level (for screen names etc.) --- 12pt bold
\titleformat{\paragraph}[block]
  {\normalfont\fontsize{12}{15}\bfseries}
  {}{0pt}{}
\titlespacing*{\paragraph}{0pt}{8pt}{4pt}

% --- Numbering ---
\renewcommand{\thechapter}{\Roman{chapter}}
\renewcommand{\thesection}{\arabic{section}}
\renewcommand{\thesubsection}{\thesection.\arabic{subsection}}
\renewcommand{\thesubsubsection}{\thesubsection.\arabic{subsubsection}}
\makeatletter
\@addtoreset{section}{chapter}
\makeatother

% --- TOC ---
\setcounter{tocdepth}{3}
\setcounter{secnumdepth}{4}
\usepackage{tocloft}

% --- Hyperlinks ---
\usepackage{hyperref}
\hypersetup{
  colorlinks=true,
  linkcolor=black,
  urlcolor=blue,
  citecolor=black
}

% --- Harvard referencing ---
\usepackage{natbib}
\bibliographystyle{agsm}

% --- Symbols ---
\usepackage{amssymb}
\usepackage{textcomp}
\usepackage{pifont}

% --- Code listings ---
\usepackage{listings}
\lstset{
  basicstyle=\ttfamily\small,
  breaklines=true,
  frame=single,
  xleftmargin=1em
}

% --- Misc ---
\usepackage{caption}
\usepackage{pdflscape}

% --- Colors for tables ---
\definecolor{tableheadpeach}{HTML}{FFE8E1}
\definecolor{placeholderblue}{HTML}{0000FF}

% --- Custom commands ---
\newcommand{\placeholder}[1]{\textcolor{placeholderblue}{\textit{[#1]}}}
\newcommand{\imgplaceholder}[1]{%
  \begin{center}
    \fbox{\begin{minipage}{0.92\textwidth}%
      \centering\vspace{5cm}\textit{\small #1}\vspace{5cm}%
    \end{minipage}}
  \end{center}
}

% ============================================================
\begin{document}

% ============================================================
%                        COVER PAGE
% ============================================================
\begin{titlepage}
\centering

\begin{tabular}{@{}c@{\hspace{10pt}}l@{}}
  % Replace with: \includegraphics[height=1.8cm]{fpt_logo.png}
  \fbox{\parbox{1.8cm}{\centering\vspace{0.3cm}{\scriptsize FPT\\Logo}\vspace{0.3cm}}} &
  \begin{tabular}[c]{@{}l@{}}
    \textbf{\fontsize{11}{13}\selectfont MINISTRY OF EDUCATION AND TRAINING}\\[2pt]
    \textbf{\fontsize{11}{13}\selectfont FPT UNIVERSITY}
  \end{tabular}
\end{tabular}

\vspace{2.5cm}

{\fontsize{26}{32}\selectfont\bfseries FPT UNIVERSITY\par}
\vspace{0.8cm}
{\fontsize{20}{24}\selectfont\textit{Capstone Project Document}\par}
\vspace{0.6cm}
{\fontsize{16}{20}\selectfont System Configuration Management (ABAP19)\par}

\vspace{2cm}

\begin{tabular}{@{}l@{\hspace{1cm}}l@{}}
  \textbf{Group Name} & GSP26 \\[8pt]
  \textbf{Group Members} &
    \begin{tabular}[t]{@{}l@{}}
      \placeholder{Lý Thúc Bảo --- Leader --- SE181969} \\
      \placeholder{Lê Minh Nghĩa --- Member --- SE181956} \\
      \placeholder{Tiêu Minh Thiện --- Member --- SE181952} \\
      \placeholder{Trần Văn Ngọc Đặng --- Member --- SE181989} \\
      \placeholder{Lê Hoàng Quý --- Member --- SE173141}
    \end{tabular} \\[8pt]
  \textbf{Supervisors}    & \begin{tabular}[t]{@{}l@{}}Nguyễn Thị Cẩm Hương \\ Nguyễn Hoàng Linh\end{tabular} \\
\end{tabular}

\vfill

{\fontsize{14}{17}\selectfont --- Hanoi, \placeholder{Month Year} ---\par}

\end{titlepage}

% ============================================================
%                     TABLE OF CONTENTS
% ============================================================
\pagenumbering{roman}
\tableofcontents
\clearpage
\listoftables
\clearpage
\listoffigures
\clearpage
\pagenumbering{arabic}
\setcounter{page}{1}

% ============================================================
%                     ACKNOWLEDGEMENT
% ============================================================
\chapter*{Acknowledgement}
\addcontentsline{toc}{chapter}{Acknowledgement}

We would like to express our sincere gratitude to FPT University for providing us with the opportunity to undertake this capstone project, which has allowed us to apply our academic knowledge to a real-world enterprise software challenge.

We extend our heartfelt thanks to our supervisors, Nguyễn Thị Cẩm Hương and Nguyễn Hoàng Linh, for the invaluable guidance and mentorship throughout the project.

Special appreciation goes to the Technical University of Munich (TUM) for granting us access to the SAP S/4HANA server (s40lp1.ucc.cit.tum.de), which served as the foundation for our development and testing environment. Without this infrastructure, the project would not have been possible.

We would also like to acknowledge the open-source and SAP developer communities whose documentation, tutorials, and forums provided essential guidance throughout the development of the RESTful ABAP Programming Model (RAP), SAP UI5, and OData V4 service integration.

Finally, we thank all members of Group GSP26 for their dedication, teamwork, and collaborative spirit that made this project possible. Each member contributed their unique skills and perspectives, resulting in a comprehensive and functional system.

% ============================================================
%                 DEFINITION AND ACRONYMS
% ============================================================
\clearpage
\chapter*{Definition and Acronyms}
\addcontentsline{toc}{chapter}{Definition and Acronyms}

The following acronyms and abbreviations are used throughout this document.

\begin{longtable}{|c|p{10cm}|}
  \hline
  \rowcolor{tableheadpeach}
  \textbf{Acronym} & \textbf{Definition} \\
  \hline
  \endfirsthead
  \hline
  \rowcolor{tableheadpeach}
  \textbf{Acronym} & \textbf{Definition} \\
  \hline
  \endhead
  \hline
  \endfoot
  \textbf{ABAP}  & Advanced Business Application Programming \\ \hline
  \textbf{API}   & Application Program Interface \\ \hline
  \textbf{BE}    & Back-End \\ \hline
  \textbf{BO}    & Business Object \\ \hline
  \textbf{BSP}   & Business Server Pages \\ \hline
  \textbf{CDS}   & Core Data Services \\ \hline
  \textbf{CRUD}  & Create, Read, Update, Delete \\ \hline
  \textbf{CSRF}  & Cross-Site Request Forgery \\ \hline
  \textbf{DCL}   & Data Control Language \\ \hline
  \textbf{DEV}   & Development Environment \\ \hline
  \textbf{ERD}   & Entity Relationship Diagram \\ \hline
  \textbf{FE}    & Front-End \\ \hline
  \textbf{FI}    & Financial Accounting \\ \hline
  \textbf{FLP}   & Fiori Launchpad \\ \hline
  \textbf{FPM}   & Flexible Programming Model \\ \hline
  \textbf{GUI}   & Graphical User Interface \\ \hline
  \textbf{KPI}   & Key Performance Indicator \\ \hline
  \textbf{LROP}  & List Report Object Page \\ \hline
  \textbf{MM}    & Materials Management \\ \hline
  \textbf{MMSS}  & Materials Management --- Safe Stock \\ \hline
  \textbf{NFR}   & Non-Functional Requirement \\ \hline
  \textbf{OData} & Open Data Protocol \\ \hline
  \textbf{PM}    & Project Manager \\ \hline
  \textbf{PRD}   & Production Environment \\ \hline
  \textbf{QAS}   & Quality Assurance System \\ \hline
  \textbf{RAP}   & RESTful ABAP Programming Model \\ \hline
  \textbf{SD}    & Sales and Distribution \\ \hline
  \textbf{SDD}   & Software Design Description \\ \hline
  \textbf{SPMP}  & Software Project Management Plan \\ \hline
  \textbf{SRS}   & Software Requirement Specification \\ \hline
  \textbf{UAT}   & User Acceptance Testing \\ \hline
  \textbf{UC}    & Use Case \\ \hline
  \textbf{UI5}   & SAP UI5 Framework \\ \hline
  \textbf{UUID}  & Universally Unique Identifier \\ \hline
\end{longtable}


% ============================================================
\clearpage
\chapter{Project Introduction}

\section{Overview}

\subsection{Project Information}

\begin{table}[H]
\centering\small
\begin{tabularx}{\textwidth}{|l|X|}
  \hline
  \rowcolor{tableheadpeach}
  \textbf{Item} & \textbf{Detail} \\ \hline
  Project Name        & System Configuration Management \\ \hline
  Project Code        & ABAP19 \\ \hline
  Group Name          & GSP26 \\ \hline
  Semester            & \placeholder{e.g., Spring 2026} \\ \hline
  University          & FPT University \\ \hline
  Technology          & SAP Fiori (UI5) + SAP RAP (ABAP) + OData V4 \\ \hline
  SAP Server          & s40lp1.ucc.cit.tum.de, Client 324 \\ \hline
  ABAP Package        & ZGSP26SAP06 \\ \hline
  GitHub Organization & \url{https://github.com/CONFIG-MNG-APP} \\ \hline
\end{tabularx}
\caption{Project Information}
\label{tab:project-info}
\end{table}

The System Configuration Management project (ABAP19) is a SAP Fiori-based enterprise application designed to centralize and govern configuration changes across multiple SAP modules --- Financial Accounting (FI), Materials Management (MM), and Sales \& Distribution (SD). The system implements a multi-tier approval workflow with environment promotion (DEV $\rightarrow$ QAS $\rightarrow$ PRD), role-based access control, comprehensive audit trailing, and real-time notifications.

The application addresses a critical gap in the SAP ecosystem: the lack of a lightweight, user-friendly tool for managing configuration changes with proper governance. While SAP provides standard table maintenance transactions (SM30/SM31), these tools offer no approval workflow, no version control, and no audit trail --- creating significant compliance and operational risks for enterprises.

\subsection{Project Team}

\begin{table}[H]
\centering\small
\begin{tabularx}{\textwidth}{|l|l|l|X|}
  \hline
  \rowcolor{tableheadpeach}
  \textbf{Name} & \textbf{Student ID} & \textbf{Role} & \textbf{Responsibilities} \\ \hline
  \placeholder{Member 1} & \placeholder{ID} & \placeholder{PM}  & \placeholder{Project planning, documentation, coordination} \\ \hline
  \placeholder{Member 2} & \placeholder{ID} & \placeholder{BE Dev}  & \placeholder{RAP backend, CDS views, behavior handlers} \\ \hline
  \placeholder{Member 3} & \placeholder{ID} & \placeholder{FE Dev}  & \placeholder{Catalog, Request, Manager apps} \\ \hline
  \placeholder{Member 4} & \placeholder{ID} & \placeholder{FE Dev}  & \placeholder{FI Limit, MM Routes, SD Price apps} \\ \hline
  \placeholder{Member 5} & \placeholder{ID} & \placeholder{FE/Test} & \placeholder{MMSS, IT Admin apps, testing} \\ \hline
\end{tabularx}
\caption{Project Team Members}
\label{tab:team}
\end{table}


\section{Product Background}

In enterprise SAP environments, system configuration parameters --- such as financial approval limits, warehouse routing rules, material safety stock levels, and sales pricing rules --- are critical to business operations. These parameters directly affect how transactions are processed, approved, and executed across the organization. A single misconfigured parameter can lead to unauthorized transactions, inventory discrepancies, or pricing errors that may have significant financial and compliance consequences.

Traditionally, SAP configuration changes are managed through standard table maintenance transactions (SM30/SM31), which present several significant challenges:

\begin{itemize}[leftmargin=2em]
  \item \textbf{No approval workflow:} Configuration changes take immediate effect without managerial review, creating compliance and audit risks. There is no mechanism for a manager to review and approve changes before they are applied to the system.
  \item \textbf{No version control:} Once a configuration is changed, the previous value is lost. There is no mechanism to compare old vs.\ new values or roll back to a previous state. This makes troubleshooting and root cause analysis extremely difficult.
  \item \textbf{No environment governance:} Changes cannot be systematically promoted from Development to Quality Assurance to Production, leading to inconsistencies across environments. Configuration drift between environments is a common source of production incidents.
  \item \textbf{No audit trail:} There is no centralized log of who changed what, when, and why. This creates significant challenges for compliance auditing (SOX, GDPR) and incident investigation.
  \item \textbf{Fragmented user experience:} Different modules (FI, MM, SD) use separate transactions with inconsistent interfaces, requiring users to learn multiple tools and workflows.
\end{itemize}

The consequences of configuration errors vary significantly across SAP modules, and the severity of each error depends on the module in which it occurs. In the Financial Accounting (FI) module, a misconfigured expense approval limit can lead to unauthorized payments being processed without managerial oversight. For example, if the automatic approval threshold is accidentally set to an excessively high value, invoices that should require manual review may be paid automatically, exposing the organization to fraud risk and financial loss. In the Materials Management (MM) module, incorrect warehouse routing rules can cause materials to be shipped to the wrong location, resulting in delayed production, increased logistics costs, and potential stockouts at the intended destination. In the Sales and Distribution (SD) module, pricing rule errors can lead to incorrect discount calculations, wrong minimum order values, or expired validity periods being applied to active customer contracts, directly affecting revenue and customer relationships.

The regulatory and compliance landscape further amplifies the importance of governed configuration management. The Sarbanes-Oxley Act (SOX) Section 404 requires publicly traded companies to establish and maintain adequate internal controls over financial reporting, which explicitly includes controls over system configuration changes that affect financial data processing. Organizations must demonstrate that configuration changes follow a documented approval process with segregation of duties --- a requirement that standard SM30/SM31 transactions cannot satisfy. Similarly, the General Data Protection Regulation (GDPR) Article 30 mandates that organizations maintain records of processing activities, including changes to system configurations that govern how personal data is handled. A centralized audit trail that captures who changed what configuration, when, and for what reason is therefore not merely a best practice but a regulatory necessity for enterprises operating in regulated industries.

Compared to traditional change management processes in enterprise IT --- where configuration changes are typically documented in spreadsheets, approved via email chains, and manually applied by basis administrators --- the approach proposed in this project offers substantial improvements. Traditional processes suffer from version conflicts when multiple stakeholders edit the same spreadsheet simultaneously, from human error when basis administrators manually type configuration values into SM30 screens, and from incomplete audit trails when email approvals are lost or archived. By embedding the entire change lifecycle (proposal, review, approval, application, promotion, and rollback) into a single integrated platform, the system eliminates these process gaps and provides a single source of truth for configuration governance.

The project customer is an enterprise organization utilizing SAP S/4HANA, represented by the academic partnership between FPT University and the Technical University of Munich (TUM), which provides access to a shared SAP S/4HANA system for development and testing. This partnership ensures that the application is developed and tested on a realistic enterprise SAP landscape.


\section{Existing Solutions}

Table~\ref{tab:existing-solutions} provides a comparative analysis of existing approaches to SAP configuration management. Each solution is evaluated against the key requirements identified in the product background.

\begin{table}[H]
\centering\small
\begin{tabularx}{\textwidth}{|l|X|X|X|}
  \hline
  \rowcolor{tableheadpeach}
  \textbf{Solution} & \textbf{Description} & \textbf{Pros} & \textbf{Cons} \\ \hline
  SM30/SM31 &
    Built-in SAP transactions for direct table editing &
    No development needed; immediate access &
    No workflow, no version control, no audit trail, poor UX \\ \hline
  SAP Solution Manager &
    Enterprise change management platform &
    Comprehensive change tracking; transport integration &
    Overly complex for config management; expensive licensing \\ \hline
  Custom ABAP Txns &
    Bespoke Z-transactions per module &
    Tailored to specific needs &
    Fragmented UX; high maintenance cost; no unified workflow \\ \hline
  Spreadsheet-based &
    Excel/CSV files tracked manually &
    Simple; no SAP development required &
    Error-prone; no real-time sync; version conflicts \\ \hline
  SAP Fiori MYS &
    Standard Fiori app for solution documentation &
    Modern UI; SAP-supported &
    Focused on documentation, not config value management \\ \hline
\end{tabularx}
\caption{Comparison of Existing Solutions}
\label{tab:existing-solutions}
\end{table}

None of the existing solutions provide a unified, Fiori-native interface that combines configuration CRUD operations with a governed approval workflow, multi-environment promotion, and comprehensive audit capabilities. This gap represents both a market opportunity and the core motivation for this project.


\section{Business Opportunity}

The gap in the SAP ecosystem for a lightweight, user-friendly configuration management tool presents a clear business opportunity.

\textbf{Market Need:} Mid-size to large enterprises running SAP S/4HANA require governed configuration changes but find SAP Solution Manager too complex and expensive for this specific use case. The growing regulatory emphasis on change management auditing (SOX compliance, GDPR data governance) further drives demand for transparent, auditable configuration workflows.

\textbf{Target Users:}
\begin{itemize}[leftmargin=2em]
  \item \textbf{Key Users} (business analysts, functional consultants) who define configuration values and propose changes through an intuitive interface
  \item \textbf{Managers} (department heads, functional leads) who approve or reject changes with full visibility into what is being modified, including side-by-side comparison of old and new values
  \item \textbf{IT Administrators} who promote approved changes across environments (DEV $\rightarrow$ QAS $\rightarrow$ PRD) with confidence in rollback capabilities
\end{itemize}

\textbf{Competitive Advantages:}
\begin{itemize}[leftmargin=2em]
  \item \textbf{SAP-native:} Built entirely on SAP technology (RAP, CDS, OData V4, Fiori), ensuring seamless integration with existing SAP landscapes without third-party dependencies
  \item \textbf{Modern UX:} SAP Fiori design language provides intuitive, responsive interfaces accessible from desktop and tablet devices
  \item \textbf{3-environment promotion:} Systematic DEV $\rightarrow$ QAS $\rightarrow$ PRD pipeline with strict ordering and rollback capability ensures production safety
  \item \textbf{Full audit trail:} Every action (create, approve, reject, promote, rollback) is logged with old/new values serialized in JSON format, enabling complete traceability
  \item \textbf{Excel import:} Bulk configuration changes via familiar spreadsheet workflows with intelligent header mapping (15--20 aliases per module)
  \item \textbf{Real-time notifications:} In-app polling (15--60 second intervals) and external email notifications keep all stakeholders informed of workflow state changes
\end{itemize}

The total addressable market for SAP configuration management tools encompasses the estimated 400,000+ organizations worldwide running SAP ERP systems, of which approximately 30,000 have migrated or are migrating to SAP S/4HANA. Among these, mid-size enterprises with 500 to 5,000 employees represent the primary market segment --- organizations large enough to require formal configuration governance but not so large that they can justify the cost and complexity of SAP Solution Manager or SAP Cloud ALM for this specific purpose. The modular architecture of the system is a key strategic advantage for market expansion: adding support for a new SAP module requires only defining a new CDS view entity with its corresponding consumption view, creating a new database table pair (request table and main table), implementing module-specific validation logic in the rule classes, and developing a new frontend application that extends the shared \texttt{ConfMngBaseController}. This pattern has been proven across four modules (FI, MM, MMSS, SD) and can be replicated for additional modules such as Plant Maintenance (PM), Quality Management (QM), or Human Capital Management (HCM) with minimal architectural changes to the existing backend infrastructure.


\section{Software Product Vision}

\textbf{Vision Statement:} To deliver a centralized, SAP Fiori-based configuration management platform that enables governed changes across multiple SAP modules with multi-tier approval, environment promotion, version control, and comprehensive audit trailing --- transforming configuration management from an ad-hoc, error-prone process into a structured, transparent, and auditable workflow.

The system empowers Key Users to propose configuration changes through an intuitive interface, Managers to review and approve changes with full visibility into what is being modified, and IT Administrators to systematically promote approved changes across Development, Quality Assurance, and Production environments with confidence in rollback capabilities.

The platform is designed to be modular and extensible, allowing additional configuration modules to be added in the future by defining new CDS views, behavior implementations, and frontend applications that follow the established patterns (ConfMngBaseController, ExcelImport, NotificationService).


\section{Project Scope \& Limitations}

\subsection{In Scope}

The following capabilities are included in the project scope:

\begin{itemize}[leftmargin=2em]
  \item \textbf{4 configuration modules:} FI Limit (expense approval limits), MM Routes (warehouse routing), MM Safe Stock (minimum inventory levels), SD Price (sales pricing rules)
  \item \textbf{3 user roles:} Key User, Manager, IT Administrator --- each with distinct permissions enforced at both frontend and backend levels
  \item \textbf{3 environments:} Development (DEV), Quality Assurance (QAS), Production (PRD) --- with strict promotion and rollback ordering
  \item \textbf{8 frontend applications:} Configuration Catalog, Change Request Management, 4 module-specific configuration editors, Manager Approval App, IT Admin Promotion App
  \item \textbf{1 backend system:} ABAP RAP with 31 database tables, 51 CDS views, 6 rule classes, 8 OData V4 services
  \item \textbf{Core workflows:} Create request $\rightarrow$ Add config lines $\rightarrow$ Submit $\rightarrow$ Approve/Reject $\rightarrow$ Apply to DEV $\rightarrow$ Promote to QAS $\rightarrow$ Promote to PRD $\rightarrow$ Rollback
  \item \textbf{Supporting features:} Excel import, real-time in-app notifications, email notifications via external service, changelog viewer, old/new value comparison, environment comparison across DEV/QAS/PRD, and audit trail with JSON diff viewer
  \item \textbf{External service:} Node.js email notification service deployed on Render.com with fire-and-forget pattern
\end{itemize}

\subsection{Out of Scope}

The following items are explicitly excluded from the project scope:

\begin{itemize}[leftmargin=2em]
  \item SAP transport management integration (CTS/CTS+)
  \item Automated production deployment pipelines
  \item Integration with SAP Solution Manager or SAP Cloud ALM
  \item Mobile-native applications (responsive web only)
  \item Configuration types beyond the 4 defined modules
  \item Multi-tenant or multi-client support (single client 324)
  \item Batch scheduling or automated configuration synchronization
\end{itemize}

\subsection{Limitations}

\begin{itemize}[leftmargin=2em]
  \item \textbf{Server dependency:} Development and testing rely on the TUM-hosted SAP S/4HANA server, which has scheduled maintenance windows and shared resource constraints
  \item \textbf{Team size:} 5-person team with limited SAP development experience at project start, requiring a significant training investment
  \item \textbf{Timeline:} Constrained to a single academic semester
  \item \textbf{Authentication:} Relies on SAP standard authentication; no custom SSO integration
  \item \textbf{Email service:} External dependency on Render.com free tier with cold-start latency of up to 30 seconds on first request
\end{itemize}


% ============================================================
%       CHAPTER II.  PROJECT MANAGEMENT PLAN
% ============================================================
\clearpage
\chapter{Project Management Plan}

\section{Overview}

\subsection{Scope \& Estimation}

Table~\ref{tab:scope-estimation} presents the scope and effort estimation for each component of the system. The complexity classification follows the Function Point Analysis approach, categorizing each component as Simple, Medium, or Complex based on the number of data elements, user interactions, and integration points.

\begin{table}[H]
\centering\small
\begin{tabularx}{\textwidth}{|c|X|c|r|}
  \hline
  \rowcolor{tableheadpeach}
  \textbf{No} & \textbf{Function / Component} & \textbf{Complexity} & \textbf{Effort (MD)} \\ \hline
  1  & Configuration Catalog App                               & Medium  &  8 \\ \hline
  2  & Change Request App (Dashboard, Landing, Object Page)    & Complex & 15 \\ \hline
  3  & FI Limit Configuration App                              & Complex & 12 \\ \hline
  4  & MM Routes Configuration App                             & Complex & 12 \\ \hline
  5  & MM Safe Stock Configuration App                         & Complex & 14 \\ \hline
  6  & SD Price Configuration App                              & Complex & 14 \\ \hline
  7  & Manager Approval App (Dashboard, ListReport, ObjectPage)& Complex & 12 \\ \hline
  8  & IT Admin App (FCL, Promote, Rollback, Audit, EnvCompare)& Complex & 18 \\ \hline
  9  & Backend RAP (CDS views, Behavior handlers, Rule classes)& Complex & 25 \\ \hline
  10 & Database Design (31 tables, seed data)                   & Medium  &  8 \\ \hline
  11 & Email Notification Service (Node.js)                     & Simple  &  4 \\ \hline
  12 & Excel Import Feature (SheetJS integration)               & Medium  &  6 \\ \hline
  13 & Notification System (polling, localStorage)              & Medium  &  5 \\ \hline
  14 & Testing (Unit, Integration, System)                      & Medium  & 10 \\ \hline
  15 & Documentation \& Deployment                              & Simple  &  7 \\ \hline
  \rowcolor{tableheadpeach}
     & \textbf{Total}                                          &         & \textbf{170} \\ \hline
\end{tabularx}
\caption{Scope \& Effort Estimation (MD = Man-Days)}
\label{tab:scope-estimation}
\end{table}


\subsection{Project Objectives}

\textbf{Quality Objectives:}
\begin{itemize}[leftmargin=2em]
  \item All 18 use cases implemented and functional
  \item Unit test coverage for critical modules (ExcelImport, Rule classes)
  \item Zero critical/high severity defects in final release
  \item Compliance with SAP Fiori design guidelines
\end{itemize}

\textbf{Time Objectives:}
\begin{itemize}[leftmargin=2em]
  \item \placeholder{Project start date} to \placeholder{Project end date}
  \item Bi-weekly sprint deliverables with demonstrable increments
\end{itemize}

\textbf{Effort Distribution:}

\begin{table}[H]
\centering\small
\begin{tabular}{|l|r|r|}
  \hline
  \rowcolor{tableheadpeach}
  \textbf{Activity} & \textbf{Percentage} & \textbf{Man-Days} \\ \hline
  Requirement Analysis                 & 10\% & 17 \\ \hline
  System Design                        & 15\% & 25 \\ \hline
  Coding (Frontend + Backend)          & 45\% & 77 \\ \hline
  Testing                              & 20\% & 34 \\ \hline
  Project Management \& Documentation  & 10\% & 17 \\ \hline
  \rowcolor{tableheadpeach}
  \textbf{Total}                       & \textbf{100\%} & \textbf{170} \\ \hline
\end{tabular}
\caption{Effort Distribution by Activity}
\label{tab:effort-distribution}
\end{table}


\subsection{Project Risks}

\begin{table}[H]
\centering\small
\begin{tabularx}{\textwidth}{|c|X|c|c|X|}
  \hline
  \rowcolor{tableheadpeach}
  \textbf{No} & \textbf{Risk} & \textbf{Prob.} & \textbf{Impact} & \textbf{Mitigation Strategy} \\ \hline
  1 & SAP server downtime or unavailability
    & Med & High & Use local mock services for FE development; schedule critical testing during off-peak hours \\ \hline
  2 & RAP/CDS learning curve for team members
    & High & Med & Conduct training in first 2 weeks; pair programming with experienced member \\ \hline
  3 & Cross-app navigation integration failures
    & Med & Med & Define and test navigation contracts early; use localStorage fallback \\ \hline
  4 & OData V4 draft handling complexity
    & High & High & Build shared ConfMngBaseController pattern; document draft choreography \\ \hline
  5 & Email service external dependency
    & Low & Low & Fire-and-forget pattern; app functions without email service \\ \hline
  6 & Data inconsistency across 3 environments
    & Med & High & Strict promote/rollback order with validation; snapshot before changes \\ \hline
  7 & Team member unavailability
    & Low & Med & Document code patterns; cross-train on critical modules \\ \hline
\end{tabularx}
\caption{Project Risk Assessment}
\label{tab:risks}
\end{table}


\section{Management Approach}

The project follows an \textbf{Agile Scrum} methodology adapted for the academic setting, with bi-weekly sprints and iterative delivery of working software. This approach was chosen because the SAP Fiori and RAP development landscape requires frequent feedback loops, and the team needed flexibility to adjust priorities as new technical challenges emerged.

\subsection{Project Process}

\imgplaceholder{INSERT: Agile/Scrum Process Flow Diagram --- Iterative cycle: Product Backlog, Sprint Planning, Sprint Execution (daily standups), Sprint Review, Sprint Retrospective, Increment. Arrow loops back to Backlog. Side note: Continuous activities --- Code reviews, CI testing, Documentation.}

\textbf{Sprint Structure (2-week cadence):}
\begin{itemize}[leftmargin=2em]
  \item \textbf{Day 1:} Sprint Planning --- select backlog items, estimate effort, assign tasks
  \item \textbf{Daily:} 15-minute standup (async via chat tool when in-person not possible)
  \item \textbf{Continuous:} Code reviews via GitHub Pull Requests
  \item \textbf{Day 13:} Sprint Review --- demonstrate working features to supervisor
  \item \textbf{Day 14:} Sprint Retrospective --- identify improvements for next sprint
\end{itemize}

\subsubsection{Adapting Scrum for SAP Development}

The team adapted the standard Scrum framework to accommodate the unique characteristics of SAP Fiori and RAP development. One of the most significant adaptations was the separation of backend and frontend sprints during the early phases of the project. Because the frontend applications depend on stable OData V4 service endpoints, the backend team worked approximately one sprint ahead of the frontend team during Sprints 1 through 4. This approach ensured that CDS views, behavior definitions, and service bindings were available and tested before the frontend developers began integrating with them. Starting from Sprint 5, once the core backend services were stable, both backend and frontend work proceeded in parallel within the same sprint, with the backend team focusing on rule class refinement, additional validation logic, and performance optimization while the frontend team implemented the remaining UI features.

Another adaptation involved the Definition of Done (DoD) criteria, which were tailored to the dual-technology nature of the project. For a backend story to be considered ``Done,'' the following criteria had to be met: (1) CDS views compile without errors in Eclipse ADT, (2) behavior definitions are syntactically correct and all handler methods are implemented, (3) OData V4 service binding is registered and accessible via the service URL, (4) ABAP Unit tests pass for any new or modified rule class methods, (5) ABAP Test Cockpit (ATC) shows no critical findings, and (6) at least one team member has reviewed the code via GitHub Pull Request. For a frontend story, the DoD required: (1) the application launches without console errors, (2) all specified UI fields and controls are present and functional, (3) value help dialogs load data correctly from the backend, (4) the draft choreography (POST, Activate, Edit, PATCH, Activate) completes without errors for all CRUD operations, (5) the application follows SAP Fiori Design Guidelines for layout, spacing, and color usage, and (6) the code has been reviewed and merged to the main branch.

Sprint velocity was tracked informally using story points estimated during Sprint Planning. The team started with an initial velocity of approximately 20 story points per sprint but found that early estimates were consistently optimistic due to the learning curve associated with SAP technologies. By Sprint 3, the team had calibrated their estimation process and stabilized at approximately 28 story points per sprint. Key factors that improved estimation accuracy included: breaking down complex stories (such as ``Implement Excel Import'') into smaller tasks with clearer scope, accounting for SAP server downtime in capacity planning, and establishing reference stories for common patterns (e.g., ``implement a new value help dialog'' was consistently estimated at 2 story points). The team also introduced a ``spike'' story type for technical investigation tasks, such as researching the OData V4 draft handling protocol or evaluating the SheetJS library for Excel parsing, which helped separate exploration work from implementation work in sprint planning.

\subsection{Quality Management}

Quality is ensured through multiple layers of verification and validation:
\begin{enumerate}[leftmargin=2em]
  \item \textbf{Code Review:} All code changes require Pull Request review by at least one team member before merging to the main branch
  \item \textbf{Unit Testing:} QUnit tests for frontend Excel import modules; ABAP Unit for backend rule classes (status machine, validator, snapshot)
  \item \textbf{Integration Testing:} OPA5 journey tests for basic app navigation flows
  \item \textbf{Manual Testing:} End-to-end workflow validation on the SAP server with real data
  \item \textbf{SAP Standards:} ABAP ATC (ABAP Test Cockpit) checks for backend code quality and adherence to SAP development guidelines
  \item \textbf{UI Guidelines:} Adherence to SAP Fiori Design Guidelines for consistent user experience across all 8 frontend applications
\end{enumerate}

\subsection{Training Plan}

Given the team's limited initial experience with SAP technologies, a structured training plan was essential. Table~\ref{tab:training} outlines the training activities conducted during the first weeks of the project.

\begin{table}[H]
\centering\small
\begin{tabularx}{\textwidth}{|X|c|X|l|}
  \hline
  \rowcolor{tableheadpeach}
  \textbf{Topic} & \textbf{Duration} & \textbf{Method} & \textbf{Target} \\ \hline
  SAP UI5 / Fiori Elements       & 1 week & openSAP courses + hands-on labs     & All FE devs \\ \hline
  RAP \& CDS View Entities       & 1 week & SAP Learning Hub + pair coding      & All BE devs \\ \hline
  OData V4 Protocol              & 3 days & Documentation study + exercises     & All devs \\ \hline
  abapGit \& ADT (Eclipse)       & 2 days & Internal workshop                   & All devs \\ \hline
  GitHub Workflow                 & 1 day  & Internal session                    & All members \\ \hline
  SAP Fiori Design Guidelines    & 2 days & SAP documentation review            & All FE devs \\ \hline
\end{tabularx}
\caption{Training Plan}
\label{tab:training}
\end{table}


\section{Project Deliverables}

Table~\ref{tab:deliverables} lists the main project deliverables, including both internal artifacts and external-facing outputs.

\begin{table}[H]
\centering\small
\begin{tabularx}{\textwidth}{|c|X|X|l|}
  \hline
  \rowcolor{tableheadpeach}
  \textbf{No} & \textbf{Deliverable} & \textbf{Description} & \textbf{Format} \\ \hline
  1 & Schedule / Task Tracking  & Sprint backlogs, burn-down charts    & GitHub / Excel \\ \hline
  2 & Project Backlog           & User stories with acceptance criteria & GitHub Issues \\ \hline
  3 & Source Codes --- Frontend  & 8 SAP UI5 applications               & GitHub repos \\ \hline
  4 & Source Codes --- Backend   & ABAP RAP package ZGSP26SAP06         & abapGit \\ \hline
  5 & Source Codes --- Mail Svc  & Node.js notification service          & GitHub repo \\ \hline
  6 & Final Report Document     & This comprehensive document           & Word / PDF \\ \hline
  7 & Test Cases Document       & Unit, Integration, System tests       & Excel (.xlsx) \\ \hline
  8 & Defects List              & Bug tracking and resolution           & GitHub Issues \\ \hline
  9 & Presentation Slides       & Final defense presentation            & PowerPoint \\ \hline
\end{tabularx}
\caption{Project Deliverables}
\label{tab:deliverables}
\end{table}


\section{Responsibility Assignments}

Table~\ref{tab:raci} shows the responsibility assignment matrix using the RACI model (R = Responsible, A = Accountable, C = Consulted, I = Informed).

\begin{table}[H]
\centering\small
\begin{tabularx}{\textwidth}{|X|c|c|c|c|c|}
  \hline
  \rowcolor{tableheadpeach}
  \textbf{Task Area} & \textbf{\placeholder{M1}} & \textbf{\placeholder{M2}} & \textbf{\placeholder{M3}} & \textbf{\placeholder{M4}} & \textbf{\placeholder{M5}} \\ \hline
  Project Management       & R,A &     &     &     &     \\ \hline
  Backend RAP Development  &     & R,A &     &     &     \\ \hline
  Catalog \& Request App   &     &     & R,A &     &     \\ \hline
  FI Limit \& MM Routes    &     &     &     & R,A &     \\ \hline
  MMSS \& SD Price App     &     &     & R   &     & A   \\ \hline
  Manager \& IT Admin App  &     & R   &     &     & A   \\ \hline
  Mail Service             &     &     &     &     & R,A \\ \hline
  Testing                  & A   &     &     &     & R   \\ \hline
  Documentation            & R   &     &     & A   &     \\ \hline
\end{tabularx}
\caption{RACI Responsibility Matrix}
\label{tab:raci}
\end{table}


\section{Project Communications}

Effective communication is critical for a distributed team working on a complex SAP project. Table~\ref{tab:communications} outlines the communication plan.

\begin{table}[H]
\centering\small
\begin{tabularx}{\textwidth}{|X|l|l|X|}
  \hline
  \rowcolor{tableheadpeach}
  \textbf{Type} & \textbf{Frequency} & \textbf{Tool} & \textbf{Participants} \\ \hline
  Daily Standup       & Daily (async)     & Zalo / Telegram   & All members \\ \hline
  Sprint Planning     & Bi-weekly         & Teams / In-person  & All members \\ \hline
  Sprint Review       & Bi-weekly         & In-person          & All + Supervisor \\ \hline
  Code Review         & Per PR            & GitHub Comments    & Dev + Reviewer \\ \hline
  Supervisor Meeting  & Weekly            & In-person / Teams  & PM + Supervisor \\ \hline
  Emergency           & As needed         & Phone / Zalo       & Relevant members \\ \hline
\end{tabularx}
\caption{Communication Plan}
\label{tab:communications}
\end{table}

\textbf{Project Interface:} The GitHub organization (\texttt{CONFIG-MNG-APP}) serves as the central hub for source code, issue tracking, and project documentation. All team members have contributor access to all 9 repositories.


\section{Configuration Management}

\subsection{Document Management}

All project documents are stored in a shared Google Drive / \placeholder{storage tool} with a structured folder hierarchy:

\begin{verbatim}
ABAP19_GSP26/
+-- 01_Requirements/
+-- 02_Design/
+-- 03_Source_Code/           (GitHub links)
+-- 04_Testing/
+-- 05_Reports/
|   +-- Report1_ProjectIntro/
|   +-- Report2_ProjectPlan/
|   +-- Report3_SRS/
|   +-- Report4_SDD/
|   +-- Report5_Testing/
|   +-- Report6_Release/
|   +-- FinalReport/
+-- 06_Presentations/
\end{verbatim}

\textbf{Versioning Convention:} \texttt{DocumentName\_vX.Y.docx} where X = major version, Y = minor revision. All document changes are tracked with a revision history table.

\subsection{Source Code Management}

\textbf{GitHub Organization:} \url{https://github.com/CONFIG-MNG-APP}

The project source code is organized into 9 repositories, as shown in Table~\ref{tab:repos}.

\begin{table}[H]
\centering\small
\begin{tabularx}{\textwidth}{|l|c|X|}
  \hline
  \rowcolor{tableheadpeach}
  \textbf{Repository} & \textbf{Branch} & \textbf{Description} \\ \hline
  config-catalog-fiori    & master & Configuration Catalog App \\ \hline
  conf-header-item        & main   & Change Request App \\ \hline
  config-fi-limit-fiori   & main   & FI Limit Configuration App \\ \hline
  config-mm-routes-fiori  & main   & MM Routes Configuration App \\ \hline
  conf-mm-safestock       & master & MM Safe Stock Configuration App \\ \hline
  config-sd-price-fiori   & main   & SD Price Configuration App \\ \hline
  config-manager-fiori    & main   & Manager Approval App \\ \hline
  conf-mng-itadmin        & master & IT Admin App \\ \hline
  configuration-mgn-rap   & main   & Backend RAP (abapGit) \\ \hline
\end{tabularx}
\caption{Source Code Repositories}
\label{tab:repos}
\end{table}

\textbf{Branching Strategy:}
\begin{itemize}[leftmargin=2em]
  \item \texttt{main} / \texttt{master} --- production-ready code
  \item \texttt{feature/<name>} --- new feature development branches
  \item \texttt{fix/<name>} --- bug fix branches
  \item Pull Request required for all merges to the main branch, with at least one reviewer approval
\end{itemize}

\subsection{Tools \& Infrastructures}

Table~\ref{tab:tools} lists all tools and infrastructure used in the project.

\begin{table}[H]
\centering\small
\begin{tabularx}{\textwidth}{|l|X|l|}
  \hline
  \rowcolor{tableheadpeach}
  \textbf{Tool} & \textbf{Purpose} & \textbf{Version / URL} \\ \hline
  SAP S/4HANA           & Backend server              & s40lp1.ucc.cit.tum.de \\ \hline
  SAP BAS / VS Code     & Frontend IDE                & Latest \\ \hline
  Eclipse ADT           & ABAP development            & Latest + ADT plugin \\ \hline
  @ui5/cli              & UI5 build \& serve          & Latest \\ \hline
  npm                   & Package management          & v9+ \\ \hline
  Node.js               & Mail service runtime        & v18+ \\ \hline
  abapGit               & ABAP version control        & Latest \\ \hline
  GitHub                & Source code hosting          & CONFIG-MNG-APP \\ \hline
  Render.com            & Mail service hosting        & Free tier \\ \hline
  SheetJS (xlsx)        & Excel file parsing          & Bundled in FE \\ \hline
  Chrome / Edge         & Testing browser             & Latest \\ \hline
\end{tabularx}
\caption{Development Tools \& Infrastructure}
\label{tab:tools}
\end{table}


% ============================================================


\clearpage
\chapter{Software Requirement Specification}

\section{Requirement Overview}

\subsection{Context Diagram}

\imgplaceholder{INSERT: System Context Diagram --- Central box: Configuration Management System. Actors: KEY USER (Create/Edit requests, Import Excel, Submit), MANAGER (Approve/Reject, View dashboards), IT ADMIN (Promote DEV to QAS to PRD, Rollback, Compare environments). External systems: SAP S/4HANA (OData V4, 8 services), Email Service on Render.com (POST /api/notify), Browser localStorage (Notification state, Nav mode).}

The System Configuration Management application is a web-based SAP Fiori system that interfaces with three categories of users and two external systems.

\textbf{Actors:}
\begin{itemize}[leftmargin=2em]
  \item \textbf{Key User:} Creates configuration change requests, adds/edits/deletes configuration lines across 4 modules (FI, MM, MMSS, SD), imports data from Excel, and submits requests for approval.
  \item \textbf{Manager:} Reviews submitted requests, compares old/new configuration values, approves or rejects requests with documented reasons.
  \item \textbf{IT Administrator:} Applies approved changes to the DEV environment, promotes configurations through QAS to PRD, performs rollbacks when needed, and monitors system state via audit trails and environment comparisons.
\end{itemize}

\textbf{External Systems:}
\begin{itemize}[leftmargin=2em]
  \item \textbf{SAP S/4HANA:} Backend system hosting the RAP-based business logic, 31 database tables, and 8 OData V4 services.
  \item \textbf{Email Notification Service:} External Node.js service deployed on Render.com that sends Gmail notifications for key workflow events (submit, approve, reject, promote, rollback).
\end{itemize}


\subsection{User Requirements}

\imgplaceholder{INSERT: Use Case Diagram (UML) --- System boundary: Configuration Management System. Actors: Key User (left), Manager (top), IT Admin (right). Key User UCs: UC-001 Create Request, UC-002 Add/Edit Lines, UC-003 Delete Lines, UC-004 Import Excel, UC-005 Save Draft, UC-006 Submit, UC-018 Resubmit. Manager UCs: UC-007 Approve, UC-008 Reject, UC-012 Changelog, UC-013 Compare. IT Admin UCs: UC-009 Apply DEV, UC-010 Promote, UC-011 Rollback, UC-014 EnvCompare, UC-015 Audit Trail. Shared: UC-016 Notifications, UC-017 Email. Relationships: UC-002 include UC-001; UC-005 include UC-002; UC-018 extend UC-006.}

\begin{table}[H]
\centering\small
\begin{tabularx}{\textwidth}{|l|l|X|}
  \hline
  \rowcolor{tableheadpeach}
  \textbf{Actor} & \textbf{Role Level} & \textbf{Description} \\ \hline
  Key User & KEY USER & Business users who define and maintain configuration values. Can create requests, edit config lines, import from Excel, save drafts, and submit for approval. Can only view their own requests. \\ \hline
  Manager & MANAGER & Department leads who review submitted requests. Can approve or reject. Have visibility into all submitted requests across modules. \\ \hline
  IT Administrator & IT ADMIN & System administrators managing environment promotion. Can apply to DEV, promote through QAS to PRD, rollback, compare environments, and view audit trails. \\ \hline
\end{tabularx}
\caption{Actor Descriptions}
\label{tab:actors}
\end{table}


\subsection{System Functionalities}

\subsubsection{Screen Flow}

\imgplaceholder{INSERT: Screen Flow Diagram --- Navigation between 8 applications: Catalog(:8080) with List and Detail views; Request(:8081) with Dashboard, CatalogLanding, ListReport, ObjectPage; Config Apps(:8082--8085) each with DynamicPage main view; Manager App with Dashboard, ListReport, ObjectPage; IT Admin App with FlexibleColumnLayout (Master list + Detail/Audit panel with Preview/Journey/Info tabs and EnvCompare/AuditDiff dialogs).}

\subsubsection{Screen Descriptions}

\begin{longtable}{|l|l|p{5.5cm}|l|}
  \hline
  \rowcolor{tableheadpeach}
  \textbf{Screen} & \textbf{App} & \textbf{Description} & \textbf{Access} \\ \hline
  \endhead
  \hline
  \endfoot
  Catalog List     & Catalog    & KPI cards (Total/Active/Inactive), module filter         & All \\ \hline
  Catalog Detail   & Catalog    & Field definitions, 3 action buttons                      & All \\ \hline
  Dashboard        & Request    & 5 KPI tiles, status-filtered request table                & All \\ \hline
  Catalog Landing  & Request    & Context-filtered list, changelog modal                    & All \\ \hline
  Request List     & Request    & ListReport, role-based filtering                         & All \\ \hline
  Request Object   & Request    & Status stepper, role-based action buttons                 & All \\ \hline
  FI Limit Main    & FI (:8084) & 8-column table, 4 value helps, Excel import              & Key User \\ \hline
  MM Routes Main   & MM (:8082) & 10-column table, SendWh$\neq$ReceiveWh validation        & Key User \\ \hline
  MMSS Main        & MMSS (:8085) & 7-column table, RAP draft choreography                 & Key User \\ \hline
  SD Price Main    & SD (:8083)   & 14-column Grid Table, date validation                  & Key User \\ \hline
  Manager Dashboard& Manager    & Animated KPIs, donut/column charts                       & Manager \\ \hline
  Manager List     & Manager    & 6-tag stats bar, selection variants                      & Manager \\ \hline
  Manager Object   & Manager    & Approve/Reject actions, Compare button                   & Manager \\ \hline
  IT Admin Inbox   & IT Admin   & FCL, request list + detail/audit panel                   & IT Admin \\ \hline
  Env Compare      & IT Admin   & DEV/QAS/PRD comparison, majority vote diff               & IT Admin \\ \hline
  Audit Diff       & IT Admin   & Old/New JSON diff with highlights                        & IT Admin \\ \hline
\end{longtable}

\subsubsection{User Roles and Authorization Matrix}

\begin{table}[H]
\centering\small
\begin{tabular}{|l|c|c|c|}
  \hline
  \rowcolor{tableheadpeach}
  \textbf{Action} & \textbf{Key User} & \textbf{Manager} & \textbf{IT Admin} \\ \hline
  View Catalog                              & $\checkmark$ & $\checkmark$ & $\checkmark$ \\ \hline
  Create Request                            & $\checkmark$ &              &              \\ \hline
  Edit / Delete Draft                       & $\checkmark$ &              &              \\ \hline
  Submit Request                            & $\checkmark$ &              &              \\ \hline
  Approve Request                           &              & $\checkmark$ &              \\ \hline
  Reject Request                            &              & $\checkmark$ &              \\ \hline
  Apply to DEV                              &              &              & $\checkmark$ \\ \hline
  Promote (DEV$\rightarrow$QAS$\rightarrow$PRD) &          &              & $\checkmark$ \\ \hline
  Rollback                                  &              &              & $\checkmark$ \\ \hline
  View Audit Trail                          &              &              & $\checkmark$ \\ \hline
  Compare Environments                      &              &              & $\checkmark$ \\ \hline
  Receive Notifications                     & $\checkmark$ & $\checkmark$ & $\checkmark$ \\ \hline
\end{tabular}
\caption{Authorization Matrix}
\label{tab:auth-matrix}
\end{table}

\subsubsection{Status Workflow Overview}

The configuration change request lifecycle is governed by a well-defined status workflow that ensures proper governance and auditability at every stage. Each status represents a distinct phase in the change management process, with specific actions available and specific roles authorized to perform those actions. The following paragraphs describe each status in detail, including the available actions, the authorized roles, and the data operations that occur during each transition.

\textbf{DRAFT} is the initial status assigned to every newly created configuration change request. When a Key User clicks ``Maintain via Request'' on the Catalog Detail page, the system creates a new request header record in \texttt{ZCONFREQH} with Status set to DRAFT. In this status, the Key User has full editing control over the request: they can add new configuration lines (ActionType=``C''), modify existing lines (ActionType=``U''), delete lines (ActionType=``X''), import data from Excel files, save drafts, and edit the request title and reason fields. The request is visible only to its creator due to the DCL-based row-level security filter (\texttt{CreatedBy = \$session.user}). Neither Managers nor IT Administrators can see or interact with DRAFT requests. The only forward transition from DRAFT is to SUBMITTED, which requires at least one configuration item to be present in the request (BR-008).

\textbf{SUBMITTED} represents the ``pending approval'' phase. When the Key User clicks ``Submit Request,'' the backend validates the status transition (DRAFT to SUBMITTED) via the \texttt{zcl\_gsp26\_rule\_status} class and verifies that the request contains at least one item. Upon successful submission, the request becomes visible to all users with the MANAGER role in the Manager App. The Key User can no longer edit the request content in this status. Two actions are available: Approve (transitions to APPROVED) and Reject (transitions to REJECTED). Both actions can only be performed by users with the MANAGER role. The system sends an email notification to relevant stakeholders when a request enters the SUBMITTED status.

\textbf{APPROVED} indicates that a Manager has reviewed and accepted the proposed configuration changes. The approval action triggers several important backend operations: the \texttt{zcl\_gsp26\_rule\_validator} class validates each configuration item for required field completeness, duplicate business key detection, and source record existence (for update and delete actions). The \texttt{zcl\_gsp26\_rule\_snapshot} class serializes the current state of the affected configuration data into JSON format and stores it in the \texttt{ZAUDITLOG} table --- this snapshot serves as the restoration point for potential future rollbacks. The approval timestamp and the approving manager's user ID are recorded in the request header. From the APPROVED status, the IT Administrator can perform three actions: Apply to DEV (transitions to ACTIVE), Promote directly (transitions to PROMOTED), or Rollback (transitions to ROLLED\_BACK).

\textbf{REJECTED} is a terminal status assigned when a Manager declines the proposed changes. The rejection action requires a mandatory reason (BR-004), which is stored in the \texttt{reject\_reason} field of the request header. The rejection reason is displayed to the Key User on the Request Object Page, providing clear feedback on why the changes were not approved. Although REJECTED is formally a terminal status, the system supports a resubmission workflow (UC-018): the Key User can re-enter edit mode on a rejected request, modify the configuration lines to address the manager's feedback, and submit the request again. This transitions the status from REJECTED back to SUBMITTED, preserving the complete audit history of the original submission, rejection, and resubmission.

\textbf{ACTIVE} means that the IT Administrator has applied the approved configuration changes to the DEV environment. During the apply action, the backend processes each configuration item according to its ActionType: items with ActionType=``C'' (Create) are inserted into the module-specific main table (e.g., \texttt{ZFILIMITCONF}), items with ActionType=``U'' (Update) modify existing records in the main table, and items with ActionType=``X'' (Delete) remove records from the main table. All operations target the DEV environment partition of the main table. An audit log entry with action type APPLY is created. From the ACTIVE status, the IT Administrator can Promote to QAS (transitions to PROMOTED) or Rollback (transitions to ROLLED\_BACK).

\textbf{PROMOTED} indicates that the configuration has been copied to a higher environment (QAS or PRD). The promote action copies the configuration data from the current environment to the next one in the pipeline (DEV to QAS, or QAS to PRD), strictly enforcing the sequential promotion order (BR-005). When the target environment is PRD (Production), the \texttt{zcl\_gsp26\_rule\_version} class queries the maximum existing version number in the target table using dynamic SQL and increments it by one, assigning the new version to the promoted records. An audit log entry with action type PROMOTE is created, and email notifications are sent to all relevant stakeholders. From the PROMOTED status, further promotion is possible (if the request has not yet reached PRD), and rollback is always available.

\textbf{ROLLED\_BACK} is the status assigned when an IT Administrator reverts the configuration changes in one or more environments. The rollback action is governed by a strict ordering rule (BR-006): if a configuration has been promoted to multiple environments, the higher environments must be rolled back first (PRD before QAS, QAS before DEV). During rollback, the \texttt{zcl\_gsp26\_rule\_snapshot} class retrieves the pre-change snapshot from the \texttt{ZAUDITLOG} table, deserializes the JSON data, and restores the previous configuration values in the target environment's main table. A new audit log entry with action type ROLLBACK is created, with the old and new data fields reversed compared to the original apply or promote entry, providing a clear record of the restoration operation.

\subsubsection{Entity Relationship Overview}

\imgplaceholder{INSERT: Entity Relationship Diagram --- ZCONFCATALOG(1)--(N)ZCONFFIELDDEF; ZCONFREQH(ReqId,EnvId)(1)--(N)ZCONFREQI(ReqItemId); ZCONFREQH(ReqId)--(N) Config Request Tables: ZFILIMITREQ, ZMMROUTECONF\_REQ, ZMMSAFESTOCK\_REQ, ZSD\_PRICE\_REQ (each with ReqId, ReqItemId, ItemId); Config Request Tables(SourceItemId)--(1) Config Main Tables: ZFILIMITCONF, ZMMROUTECONF, ZMMSAFESTOCK, ZSD\_PRICE\_CONF (each with ItemId); ZAUDITLOG logged by rule\_writer; ZUSERROLE for authorization; ZENVDEF (DEV,QAS,PRD); 11 Value Help tables.}


% ------------------------------------------------------------
\section{Functional Specifications}

\subsection{Configuration CRUD}

% -- UC-001 --
\subsubsection{UC-001: Create Configuration Request}

\begin{longtable}{|>{\bfseries}p{2.8cm}|p{12cm}|}
  \hline
  \rowcolor{tableheadpeach}
  \textbf{Field} & \textbf{Description} \\ \hline
  \endhead \hline \endfoot
  UC Code       & UC-001 \\ \hline
  UC Name       & Create Configuration Request \\ \hline
  Actor         & Key User \\ \hline
  Precondition  & User is logged in; a configuration catalog entry exists and is active \\ \hline
  Main Flow     & 1.~User navigates to Catalog Detail and clicks ``Maintain''.
                  2.~System fetches CSRF token.
                  3.~System calls \texttt{createRequest} action on ZC\_CONF\_REQ\_H with ConfId, ModuleId, ConfName, TargetCds, ActionType=``U'', TargetEnvId=``DEV''.
                  4.~System queries the new request (\$top=1, ordered by CreatedAt desc) to get ReqId.
                  5.~System navigates to module-specific config app with ReqId and context params.
                  6.~Config app loads in edit mode with current DEV data. \\ \hline
  Alt Flow      & A1: User clicks ``Create Request'' from config app --- system prompts for Reason first. \\ \hline
  Postcondition & New request with Status=DRAFT and one item linked to the configuration. \\ \hline
  Business Rules& BR-007 (initial status must be DRAFT) \\ \hline
\end{longtable}

% -- UC-002 --
\subsubsection{UC-002: Add/Edit Configuration Lines}

\begin{longtable}{|>{\bfseries}p{2.8cm}|p{12cm}|}
  \hline
  \rowcolor{tableheadpeach}
  \textbf{Field} & \textbf{Description} \\ \hline
  \endhead \hline \endfoot
  UC Code       & UC-002 \\ \hline
  UC Name       & Add/Edit Configuration Lines \\ \hline
  Actor         & Key User \\ \hline
  Precondition  & Request in DRAFT status; user in edit mode \\ \hline
  Main Flow     & 1.~User clicks ``Add Line'' to insert blank row.
                  2.~User fills required fields (varies by module).
                  3.~User optionally edits existing rows (tracked as ActionType=``U'').
                  4.~System validates client-side (required fields, business rules).
                  5.~User clicks ``Save Draft''.
                  6.~For new rows: POST draft $\rightarrow$ Activate.
                  7.~For modified rows: Edit $\rightarrow$ PATCH $\rightarrow$ Activate.
                  8.~System updates request header via \texttt{updateReason} action.
                  9.~System reloads table with overlay (main + request changes). \\ \hline
  Alt Flow      & A1: Validation fails --- red borders on fields, error messages. Save blocked. \\ \hline
  Postcondition & Changes saved as active entities linked to request. \\ \hline
  Business Rules& BR-001 (SendWh$\neq$ReceiveWh), BR-002 (ValidFrom$\leq$ValidTo), BR-003 (numeric $\geq$0) \\ \hline
\end{longtable}

% -- UC-003 --
\subsubsection{UC-003: Delete Configuration Lines}

\begin{longtable}{|>{\bfseries}p{2.8cm}|p{12cm}|}
  \hline
  \rowcolor{tableheadpeach}
  \textbf{Field} & \textbf{Description} \\ \hline
  \endhead \hline \endfoot
  UC Code       & UC-003 \\ \hline
  UC Name       & Delete Configuration Lines \\ \hline
  Actor         & Key User \\ \hline
  Precondition  & Request in DRAFT status; at least one row exists \\ \hline
  Main Flow     & 1.~User selects rows (multi-select).
                  2.~Clicks ``Delete'' button.
                  3.~Confirmation dialog.
                  4.~New rows (never saved): removed immediately.
                  5.~Existing rows: marked ActionType=``X'', \_state=``deleted'' (red highlight).
                  6.~On save: deleted rows persisted via POST/PATCH with ActionType=``X''. \\ \hline
  Postcondition & Selected rows marked for deletion; changes take effect after save. \\ \hline
\end{longtable}

% -- UC-004 --
\subsubsection{UC-004: Import from Excel}

\begin{longtable}{|>{\bfseries}p{2.8cm}|p{12cm}|}
  \hline
  \rowcolor{tableheadpeach}
  \textbf{Field} & \textbf{Description} \\ \hline
  \endhead \hline \endfoot
  UC Code       & UC-004 \\ \hline
  UC Name       & Import Configuration Lines from Excel \\ \hline
  Actor         & Key User \\ \hline
  Precondition  & Request in DRAFT status; edit mode active \\ \hline
  Main Flow     & 1.~User clicks ``Import from Excel''.
                  2.~File picker dialog (.xls, .xlsx, .csv).
                  3.~System lazy-loads SheetJS (xlsx.min.js).
                  4.~Reads workbook, extracts first sheet.
                  5.~Maps headers case-insensitively (15--20 aliases per module).
                  6.~Transforms rows; parses numbers and booleans.
                  7.~If $>$500 rows: confirmation dialog.
                  8.~Prepends rows with ActionType=``C'', \_state=``new''.
                  9.~Toast: ``X rows imported (Y skipped)''. \\ \hline
  Alt Flow      & A1: No recognizable headers $\rightarrow$ error. A2: All empty $\rightarrow$ warning. \\ \hline
  Postcondition & Imported rows appear in table ready for review and save. \\ \hline
  Business Rules& Module-specific header alias mappings \\ \hline
\end{longtable}

% -- UC-005 --
\subsubsection{UC-005: Save Draft}

\begin{longtable}{|>{\bfseries}p{2.8cm}|p{12cm}|}
  \hline
  \rowcolor{tableheadpeach}
  \textbf{Field} & \textbf{Description} \\ \hline
  \endhead \hline \endfoot
  UC Code       & UC-005 \\ \hline
  UC Name       & Save Draft Changes \\ \hline
  Actor         & Key User \\ \hline
  Precondition  & Request in DRAFT; changes exist \\ \hline
  Main Flow     & 1.~Validate all rows via \_validateRows().
                  2.~Fetch CSRF token.
                  3.~Execute OData operations per row (POST/Activate or Edit/PATCH/Activate).
                  4.~Patch request header (title, reason).
                  5.~Set editMode=false.
                  6.~Toast: ``X change(s) saved''.
                  7.~Reload table with overlay. \\ \hline
  Alt Flow      & A1: Validation error $\rightarrow$ MessageBox with all issues; save blocked. \\ \hline
  Postcondition & All changes persisted; edit mode disabled. \\ \hline
\end{longtable}


\subsection{Approval Workflow}

% -- UC-006 --
\subsubsection{UC-006: Submit Request}

\begin{longtable}{|>{\bfseries}p{2.8cm}|p{12cm}|}
  \hline
  \rowcolor{tableheadpeach}
  \textbf{Field} & \textbf{Description} \\ \hline
  \endhead \hline \endfoot
  UC Code       & UC-006 \\ \hline
  UC Name       & Submit Request for Approval \\ \hline
  Actor         & Key User \\ \hline
  Precondition  & Request in DRAFT; $\geq$1 item \\ \hline
  Main Flow     & 1.~Click ``Submit Request''.
                  2.~Validate ReqId, Status=DRAFT, item count.
                  3.~Confirmation dialog.
                  4.~Patch header via \texttt{updateReason}.
                  5.~Call \texttt{submit} bound action.
                  6.~Send email (fire-and-forget).
                  7.~Update model: Status=``SUBMITTED''.
                  8.~Success dialog; navigate back. \\ \hline
  Postcondition & Status = SUBMITTED; email sent. \\ \hline
  Business Rules& BR-007 (DRAFT$\rightarrow$SUBMITTED), BR-008 ($\geq$1 item) \\ \hline
\end{longtable}

% -- UC-007 --
\subsubsection{UC-007: Approve Request}

\begin{longtable}{|>{\bfseries}p{2.8cm}|p{12cm}|}
  \hline
  \rowcolor{tableheadpeach}
  \textbf{Field} & \textbf{Description} \\ \hline
  \endhead \hline \endfoot
  UC Code       & UC-007 \\ \hline
  UC Name       & Approve Request \\ \hline
  Actor         & Manager \\ \hline
  Precondition  & Request SUBMITTED; user has MANAGER role \\ \hline
  Main Flow     & 1.~Manager opens request.
                  2.~Reviews details and changes (Compare button).
                  3.~Clicks ``Approve''.
                  4.~Backend validates transition SUBMITTED$\rightarrow$APPROVED.
                  5.~Validates each item (required fields, no duplicates, source exists).
                  6.~Creates approval snapshot for rollback.
                  7.~Creates audit log entries.
                  8.~Updates Status=APPROVED, ApprovedBy, ApprovedAt.
                  9.~Push notification + email. \\ \hline
  Alt Flow      & A1: Validation fails $\rightarrow$ error returned. \\ \hline
  Postcondition & Status = APPROVED; snapshot created. \\ \hline
  Business Rules& BR-007 (SUBMITTED$\rightarrow$APPROVED), BR-009 (per-module validation) \\ \hline
\end{longtable}

% -- UC-008 --
\subsubsection{UC-008: Reject Request}

\begin{longtable}{|>{\bfseries}p{2.8cm}|p{12cm}|}
  \hline
  \rowcolor{tableheadpeach}
  \textbf{Field} & \textbf{Description} \\ \hline
  \endhead \hline \endfoot
  UC Code       & UC-008 \\ \hline
  UC Name       & Reject Request \\ \hline
  Actor         & Manager \\ \hline
  Precondition  & Request SUBMITTED; MANAGER role \\ \hline
  Main Flow     & 1.~Click ``Reject''.
                  2.~Prompt for rejection reason (mandatory).
                  3.~Backend validates transition and reason.
                  4.~Updates Status=REJECTED, RejectReason, RejectedBy.
                  5.~Audit log entry.
                  6.~Email to request creator. \\ \hline
  Postcondition & Status = REJECTED (terminal). Creator notified. \\ \hline
  Business Rules& BR-004 (reason required), BR-007, BR-010 (terminal) \\ \hline
\end{longtable}


\subsection{Environment Promotion}

% -- UC-009 --
\subsubsection{UC-009: Apply to DEV}

\begin{longtable}{|>{\bfseries}p{2.8cm}|p{12cm}|}
  \hline
  \rowcolor{tableheadpeach}
  \textbf{Field} & \textbf{Description} \\ \hline
  \endhead \hline \endfoot
  UC Code       & UC-009 \\ \hline
  UC Name       & Apply Approved Changes to DEV \\ \hline
  Actor         & IT Administrator \\ \hline
  Precondition  & Request APPROVED \\ \hline
  Main Flow     & 1.~Select approved request; click ``Apply''.
                  2.~Backend processes each item: CREATE$\rightarrow$insert, UPDATE$\rightarrow$modify, DELETE$\rightarrow$remove.
                  3.~Updates line status to ACTIVE.
                  4.~Audit log (action=APPLY).
                  5.~Status = ACTIVE. \\ \hline
  Postcondition & Config changes live in DEV. \\ \hline
\end{longtable}

% -- UC-010 --
\subsubsection{UC-010: Promote Configuration}

\begin{longtable}{|>{\bfseries}p{2.8cm}|p{12cm}|}
  \hline
  \rowcolor{tableheadpeach}
  \textbf{Field} & \textbf{Description} \\ \hline
  \endhead \hline \endfoot
  UC Code       & UC-010 \\ \hline
  UC Name       & Promote (DEV$\rightarrow$QAS$\rightarrow$PRD) \\ \hline
  Actor         & IT Administrator \\ \hline
  Precondition  & Request APPROVED or ACTIVE \\ \hline
  Main Flow     & 1.~Click ``Promote''.
                  2.~Determine next env (DEV$\rightarrow$QAS or QAS$\rightarrow$PRD).
                  3.~Confirmation dialog.
                  4.~Backend copies config data to next env per module.
                  5.~If target = PRD: increment version.
                  6.~Audit log (action=PROMOTE).
                  7.~Update EnvId = nextEnv.
                  8.~Email notification. \\ \hline
  Postcondition & Config data copied to next environment. \\ \hline
  Business Rules& BR-005 (strict DEV$\rightarrow$QAS$\rightarrow$PRD, no skip) \\ \hline
\end{longtable}

% -- UC-011 --
\subsubsection{UC-011: Rollback Configuration}

\begin{longtable}{|>{\bfseries}p{2.8cm}|p{12cm}|}
  \hline
  \rowcolor{tableheadpeach}
  \textbf{Field} & \textbf{Description} \\ \hline
  \endhead \hline \endfoot
  UC Code       & UC-011 \\ \hline
  UC Name       & Rollback Configuration \\ \hline
  Actor         & IT Administrator \\ \hline
  Precondition  & Request APPROVED or ACTIVE \\ \hline
  Main Flow     & 1.~Click ``Rollback'' on journey UI or table.
                  2.~Check rollback order: PRD before QAS, QAS before DEV.
                  3.~Confirmation dialog.
                  4.~Backend calls \texttt{restore\_from\_snapshot}: deserializes old\_data JSON, restores previous values.
                  5.~New audit log (action=ROLLBACK, reversed old/new).
                  6.~Status = ROLLED\_BACK.
                  7.~Email notification. \\ \hline
  Alt Flow      & A1: Order violated $\rightarrow$ error with blocking env name. \\ \hline
  Postcondition & Config restored to pre-change state. \\ \hline
  Business Rules& BR-006 (strict PRD$\rightarrow$QAS$\rightarrow$DEV) \\ \hline
\end{longtable}


\subsection{Supporting Features}

The following use cases describe auxiliary system features that support the core workflows.

% -- UC-012 --
\subsubsection{UC-012: View Changelog}

\begin{longtable}{|>{\bfseries}p{2.8cm}|p{12cm}|}
  \hline
  \rowcolor{tableheadpeach}
  \textbf{Field} & \textbf{Description} \\ \hline
  \endhead \hline \endfoot
  UC Code       & UC-012 \\ \hline
  UC Name       & View Changelog \\ \hline
  Actor         & All roles (Key User, Manager, IT Administrator) \\ \hline
  Precondition  & User is logged in; at least one request exists for the selected configuration \\ \hline
  Main Flow     & 1.~User clicks ``Show Changelog'' on a request row in the Catalog Landing page.
                  2.~System fetches config items from the module-specific OData service filtered by ReqId.
                  3.~System normalizes rows by detecting ActionType: C=ADDED, U=MODIFIED, X=DELETED.
                  4.~System maps old/new values for key fields and diff fields per module.
                  5.~A dynamic table is built with columns: Action badge, key fields, and diff fields (with old$\rightarrow$new arrows for changes).
                  6.~Rows are color-coded: green (added), orange (modified), red (deleted).
                  7.~Table is displayed inside a responsive dialog. \\ \hline
  Alt Flow      & A1: No items found for ReqId --- dialog shows ``No changes recorded'' message. \\ \hline
  Postcondition & Changelog dialog displayed; no data modified. \\ \hline
  Business Rules& Color coding follows SAP Fiori semantic colors (Positive, Critical, Negative). \\ \hline
\end{longtable}

% -- UC-013 --
\subsubsection{UC-013: Compare Old/New Values}

\begin{longtable}{|>{\bfseries}p{2.8cm}|p{12cm}|}
  \hline
  \rowcolor{tableheadpeach}
  \textbf{Field} & \textbf{Description} \\ \hline
  \endhead \hline \endfoot
  UC Code       & UC-013 \\ \hline
  UC Name       & Compare Old/New Values \\ \hline
  Actor         & Manager, IT Administrator \\ \hline
  Precondition  & Request has at least one item with ActionType = U (Update) or X (Delete) \\ \hline
  Main Flow     & 1.~User clicks ``Compare'' button on a request item row.
                  2.~System reads OldValue and NewValue fields from the request item entity.
                  3.~System parses the field values and identifies changed fields.
                  4.~A side-by-side comparison dialog is displayed.
                  5.~Changed fields are highlighted: red background for old values, green background for new values.
                  6.~Unchanged fields are displayed in normal style for context. \\ \hline
  Alt Flow      & A1: ActionType = C (Create) --- no old values exist; dialog shows new values only with ``New Record'' label.
                  A2: ActionType = X (Delete) --- dialog shows old values only with ``Deleted Record'' label. \\ \hline
  Postcondition & Comparison dialog displayed; no data modified. \\ \hline
  Business Rules& Comparison is read-only; users cannot edit values from this dialog. \\ \hline
\end{longtable}

% -- UC-014 --
\subsubsection{UC-014: Environment Comparison}

\begin{longtable}{|>{\bfseries}p{2.8cm}|p{12cm}|}
  \hline
  \rowcolor{tableheadpeach}
  \textbf{Field} & \textbf{Description} \\ \hline
  \endhead \hline \endfoot
  UC Code       & UC-014 \\ \hline
  UC Name       & Environment Comparison \\ \hline
  Actor         & IT Administrator \\ \hline
  Precondition  & User has IT ADMIN role; at least one module has configuration data in one or more environments \\ \hline
  Main Flow     & 1.~IT Admin opens Env Compare dialog and selects a module from the dropdown.
                  2.~System fetches config data from all 3 environments (DEV/QAS/PRD) in parallel.
                  3.~System builds key$\rightarrow$data maps per environment and unions all keys.
                  4.~For each field, the system calculates the majority value (mode) across environments.
                  5.~States are assigned: Success (matches majority), Warning (differs from majority), Error (missing in this environment).
                  6.~Result is displayed in a table with columns: Record Key, Field, DEV value, QAS value, PRD value, Sync status.
                  7.~Version numbers are also compared per record. \\ \hline
  Alt Flow      & A1: User toggles ``Diffs Only'' --- table filters to show only records with at least one difference.
                  A2: No data in any environment --- table is empty with informational message. \\ \hline
  Postcondition & Comparison table displayed; no data modified. \\ \hline
  Business Rules& Majority vote uses mode calculation: if 2 of 3 environments agree, the third is marked as divergent. If all 3 differ, all are marked Warning. Growing table threshold is 200 records. \\ \hline
\end{longtable}

% -- UC-015 --
\subsubsection{UC-015: View Audit Trail}

\begin{longtable}{|>{\bfseries}p{2.8cm}|p{12cm}|}
  \hline
  \rowcolor{tableheadpeach}
  \textbf{Field} & \textbf{Description} \\ \hline
  \endhead \hline \endfoot
  UC Code       & UC-015 \\ \hline
  UC Name       & View Audit Trail \\ \hline
  Actor         & IT Administrator \\ \hline
  Precondition  & User has IT ADMIN role; audit log entries exist in ZAUDITLOG \\ \hline
  Main Flow     & 1.~IT Admin views the Audit Trail panel (default mid-column view in FCL layout).
                  2.~System fetches ZAUDITLOG entries, joined with request titles from ZC\_CONF\_REQ\_H.
                  3.~Entries display: action type (APPROVE/PROMOTE/ROLLBACK), module, environment, user, and timestamp.
                  4.~Entries are sorted by timestamp descending (most recent first).
                  5.~User clicks on an entry to open the AuditDiffDialog.
                  6.~AuditDiffDialog parses old/new JSON data into a two-column table.
                  7.~Changed fields are highlighted: red for old values, green for new values. \\ \hline
  Alt Flow      & A1: No audit entries exist --- panel shows ``No audit entries found'' placeholder.
                  A2: JSON parsing fails --- dialog shows raw JSON text as fallback. \\ \hline
  Postcondition & Audit trail displayed; data is immutable and read-only. \\ \hline
  Business Rules& Audit log entries are immutable (append-only). JSON serialization uses \texttt{/ui2/cl\_json}. \\ \hline
\end{longtable}

% -- UC-016 --
\subsubsection{UC-016: In-App Notifications}

\begin{longtable}{|>{\bfseries}p{2.8cm}|p{12cm}|}
  \hline
  \rowcolor{tableheadpeach}
  \textbf{Field} & \textbf{Description} \\ \hline
  \endhead \hline \endfoot
  UC Code       & UC-016 \\ \hline
  UC Name       & In-App Notifications \\ \hline
  Actor         & All roles (Key User, Manager, IT Administrator) \\ \hline
  Precondition  & User is logged in; NotificationService is initialized on app start \\ \hline
  Main Flow     & 1.~NotificationService polls ZC\_CONF\_REQ\_H at configured intervals (15--60 seconds depending on the app).
                  2.~Service compares current status against last-seen status (stored in the \_mSeen map).
                  3.~On status change (excluding DRAFT$\rightarrow$SUBMITTED, which is the user's own action), a notification is created.
                  4.~Notification includes: title, message, status icon, and timestamp.
                  5.~Bell icon in the app header shows an unread count badge.
                  6.~User clicks the bell to open a popover with the notification list.
                  7.~Clicking a notification marks it as read and navigates to the relevant request. \\ \hline
  Alt Flow      & A1: OData fetch fails --- error is silently caught; polling continues on next interval.
                  A2: localStorage is full --- oldest read notifications are pruned first. \\ \hline
  Postcondition & Notifications displayed in bell popover; state persisted in localStorage. \\ \hline
  Business Rules& TTL for read notifications: 7 days. Maximum stored notifications: 50. Polling intervals: Catalog 60s, Request 30s, Config apps 15s, Manager 30s, IT Admin 15s. \\ \hline
\end{longtable}

% -- UC-017 --
\subsubsection{UC-017: Email Notifications}

\begin{longtable}{|>{\bfseries}p{2.8cm}|p{12cm}|}
  \hline
  \rowcolor{tableheadpeach}
  \textbf{Field} & \textbf{Description} \\ \hline
  \endhead \hline \endfoot
  UC Code       & UC-017 \\ \hline
  UC Name       & Email Notifications \\ \hline
  Actor         & System (automated) \\ \hline
  Precondition  & Mail service is deployed and accessible at the configured URL \\ \hline
  Main Flow     & 1.~On key workflow events (submit, approve, reject, promote, rollback), the frontend sends a fire-and-forget POST request to the external mail service.
                  2.~Target URL: \texttt{https://abap19-mail-206.onrender.com/api/notify}.
                  3.~Payload includes: event type, ReqId, request title, module, environment, and the triggering user.
                  4.~Mail service validates the API key from the \texttt{x-api-key} header.
                  5.~Mail service builds an HTML email template with event-specific formatting.
                  6.~Email is sent via nodemailer to Gmail API.
                  7.~Relevant stakeholders receive the formatted notification email. \\ \hline
  Alt Flow      & A1: Mail service unavailable (cold-start delay up to 30s) --- error is silently ignored; application continues normally.
                  A2: Invalid API key --- mail service returns 401; frontend ignores the error. \\ \hline
  Postcondition & Email sent to stakeholders (best-effort); application state unaffected by delivery success or failure. \\ \hline
  Business Rules& Fire-and-forget pattern: frontend does not await response. All errors silently caught. App functions fully without email service. \\ \hline
\end{longtable}

% -- UC-018 --
\subsubsection{UC-018: Resubmit Rejected Request}

\begin{longtable}{|>{\bfseries}p{2.8cm}|p{12cm}|}
  \hline
  \rowcolor{tableheadpeach}
  \textbf{Field} & \textbf{Description} \\ \hline
  \endhead \hline \endfoot
  UC Code       & UC-018 \\ \hline
  UC Name       & Resubmit Rejected Request \\ \hline
  Actor         & Key User \\ \hline
  Precondition  & Request in REJECTED status; user is the original creator \\ \hline
  Main Flow     & 1.~Key User opens the rejected request from the Request list.
                  2.~User reviews the rejection reason displayed in the request header.
                  3.~User clicks ``Edit'' to enter edit mode on the configuration lines.
                  4.~User modifies configuration lines to address the manager's feedback.
                  5.~User saves the draft changes (UC-005).
                  6.~User clicks ``Submit'' to resubmit the request for approval.
                  7.~Status transitions from REJECTED $\rightarrow$ SUBMITTED.
                  8.~Email notification sent to managers. \\ \hline
  Alt Flow      & A1: User decides not to resubmit --- request remains in REJECTED status.
                  A2: User creates a brand new request instead of resubmitting. \\ \hline
  Postcondition & Request status changed to SUBMITTED; full audit chain preserved from original submission through rejection and resubmission. \\ \hline
  Business Rules& BR-007 allows REJECTED $\rightarrow$ SUBMITTED transition. The complete audit history (original submission, rejection reason, resubmission) is maintained for traceability. \\ \hline
\end{longtable}


% ------------------------------------------------------------
\section{UI Requirements}

\subsection{Catalog App}

\paragraph{a. Catalog List Screen}

\imgplaceholder{INSERT: Screenshot --- Catalog List page. DynamicPage with pinned header. Header contains: Module Info cards (MM/SD/FI with counts), KPI cards (Total blue, Active green, Inactive grey). Content: SegmentedButton filter (ALL/MM/SD/FI), SearchField. Table columns: Module, Config Name + ID, Type, Description, Target CDS (monospace), Status. Growing=true, 20 rows per batch.}

\noindent\textbf{Related Use Cases:} UC-001 (entry point for Maintain action).

\begin{table}[H]
\centering\small
\begin{tabularx}{\textwidth}{|l|l|X|}
  \hline
  \rowcolor{tableheadpeach}
  \textbf{Component} & \textbf{Type} & \textbf{Description} \\ \hline
  Module Info Cards & ObjectStatus & Count per module (MM=Information, SD=Success, FI=Warning) \\ \hline
  KPI Cards         & Custom VBox  & Total (blue gradient), Active (green), Inactive (grey) \\ \hline
  Module Filter     & SegmentedButton & Segments: ALL, MM, SD, FI --- client-side filter \\ \hline
  Search Field      & SearchField  & Searches: ConfName, ModuleId, TargetCds, Description \\ \hline
  Notification Bell & Button+Badge & Unread notification count, opens NotificationPopover \\ \hline
\end{tabularx}
\caption{Catalog List --- UI Components}
\label{tab:ui-catalog-list}
\end{table}

\paragraph{b. Catalog Detail Screen}

\imgplaceholder{INSERT: Screenshot --- Catalog Detail page. Header: ConfName title, module+status ObjectStatus. 3 action buttons: View Current Config, View Requests, Maintain via Request. Info cards: Config Details, Description, Admin Data. Field Definitions table: FieldName, FieldLabel, DataType, ValueHelpType, IsRequired (checkmark).}


\subsection{Request App}

\paragraph{a. Dashboard Screen}

\imgplaceholder{INSERT: Screenshot --- Request Dashboard. 5 KPI tiles (Total, Submitted, Approved, Rejected, Draft) with NumericContent. Status filter SegmentedButton. Request table: Title, Module, Env, Status, Created, Actions. Notification bell.}

\paragraph{b. Catalog Landing Screen}

\imgplaceholder{INSERT: Screenshot --- Catalog Landing. Config context header (Module, Env, TargetCds). KPI cards. Request table filtered by ConfId. Show Changelog button per row opens modal with dynamic columns: Action badge, key fields, diff fields (old to new arrows).}

\paragraph{c. Request Object Page}

\imgplaceholder{INSERT: Screenshot --- Request Object Page. ObjectPageLayout. Header: ReqTitle, Status badge (inverted). Buttons: Edit, Delete, Approve, Reject, Rollback, Open Config (visible per role+status). Sections: Status Stepper (circles + lines for Draft, Submitted, Approved, Active), Request Info, Configuration, Admin Info.}


\subsection{Configuration Apps (FI / MM / MMSS / SD)}

\paragraph{a. FI Limit Configuration Screen}

\imgplaceholder{INSERT: Screenshot --- FI Limit. DynamicPage. Header cards: Request (ID+Status), Configuration (Module+Name), Request Details (Title input, Reason textarea). Summary stats: Total, Create (green), Update (blue), Delete (red). Toolbar: Add Line, Import Excel, Delete, Save/Cancel, Refresh. 8-column table: Action, Env, ExpenseType, GlAccount, AutoApprLim, Currency, Version, ChangeNote. Row highlights: green=new, blue=modified, red=deleted. Value helps for ExpenseType, GlAccount, Currency. Validation: red borders for missing required fields. Duplicate key check: ExpenseType+GlAccount.}

\begin{table}[H]
\centering\small
\begin{tabularx}{\textwidth}{|l|l|l|c|X|}
  \hline
  \rowcolor{tableheadpeach}
  \textbf{Field} & \textbf{Type} & \textbf{Width} & \textbf{Req.} & \textbf{Description} \\ \hline
  Action & ObjectStatus & 7rem & -- & Action type badge (C=Create, U=Update, X=Delete) \\ \hline
  Environment & ObjectStatus & 5.5rem & -- & Target environment (DEV/QAS/PRD) \\ \hline
  Expense Type & Input/Text & 12rem & Yes & Expense category code (value help available) \\ \hline
  GL Account & Input/Text & 10rem & Yes & General Ledger account number (value help) \\ \hline
  Approval Limit & Input (Number) & 11rem & Yes & Maximum auto-approval amount (must be >= 0) \\ \hline
  Currency & Input/Text & 9rem & Yes & Currency code (VND/USD/EUR, value help) \\ \hline
  Version & Text & 4.5rem & -- & Version number (read-only) \\ \hline
  Change Note & Input/Text & flexible & -- & Description of change (max 255 chars) \\ \hline
\end{tabularx}
\caption{FI Limit --- UI Field Descriptions}
\label{tab:ui-fi-fields}
\end{table}

The FI Limit configuration module governs expense approval thresholds within the Financial Accounting domain. In a typical enterprise scenario, different categories of expenses (such as travel, office supplies, or consulting fees) require different levels of managerial oversight depending on the amount involved. The \texttt{auto\_appr\_lim} field defines the monetary threshold below which expenses of a given type and GL account combination can be automatically approved without manual intervention. This configuration directly affects the accounts payable process: if the threshold is set incorrectly, either too many invoices bypass approval (creating fraud risk) or too many trivial expenses queue up for manual review (creating operational bottlenecks).

Key validation rules enforced by the FI Limit configuration app include: (1) ExpenseType and GlAccount are mandatory fields, ensuring that every approval limit is associated with a specific expense category and ledger account; (2) the AutoApprLim value must be a non-negative number (BR-003), preventing the entry of invalid negative thresholds; (3) the Currency field must contain a valid currency code selected from the value help dialog, which queries the \texttt{ZI\_VH\_CURRENCY} CDS view; and (4) the composite business key (ExpenseType + GlAccount) must be unique within the same request, preventing duplicate entries that would create ambiguity in the approval logic (BR-013). The backend additionally validates during the approve action that the source record exists for update and delete operations, ensuring referential integrity.

The Excel import feature for the FI Limit module accepts spreadsheet files with the following column mappings (case-insensitive, supporting 15+ aliases per field): ``Expense Type'' or ``ExpType'' or ``expense\_type'' maps to the ExpenseType field; ``GL Account'' or ``GlAccount'' or ``gl\_account'' or ``Account'' maps to GlAccount; ``Approval Limit'' or ``AutoApprLim'' or ``Limit'' or ``Amount'' maps to AutoApprLim; ``Currency'' or ``Curr'' or ``CurrencyCode'' maps to Currency; and ``Change Note'' or ``Note'' or ``Description'' maps to ChangeNote. The import parser uses the \texttt{ExcelImport.js} module, which normalizes headers by converting to lowercase and stripping whitespace before matching against the alias dictionary. Imported rows are assigned ActionType=``C'' and appear with green highlighting in the table, ready for review before saving.

The draft choreography for the FI Limit module follows the standard pattern shared across all four configuration apps. When a user clicks ``Save Draft,'' the controller iterates over all changed rows. For each new row (state=``new''), the system performs a POST request to create a draft entity with \texttt{IsActiveEntity=false}, followed by an immediate Activate action to convert the draft into an active entity. For each modified row (state=``modified''), the system first calls the Edit action on the existing active entity to create a draft copy, then PATCHes the draft with the changed field values, and finally Activates the draft. This three-step choreography (Edit, PATCH, Activate) ensures that the OData V4 draft handling protocol is followed correctly while maintaining explicit control over the save process.

\paragraph{b. MM Routes Configuration Screen}

\imgplaceholder{INSERT: Screenshot --- MM Routes. Same layout as FI. 10-column table: Action, Env, Plant, SendWh, ReceiveWh, TransMode, IsAllowed, Inspector, Version, ChangeNote. Special validation: SendWh must differ from ReceiveWh (BR-001). Boolean IsAllowed as checkbox in edit mode. Inspector value help: composite filter (UserId OR Fullname). Duplicate key check: PlantId+SendWh+ReceiveWh.}

\begin{table}[H]
\centering\small
\begin{tabularx}{\textwidth}{|l|l|l|c|X|}
  \hline
  \rowcolor{tableheadpeach}
  \textbf{Field} & \textbf{Type} & \textbf{Width} & \textbf{Req.} & \textbf{Description} \\ \hline
  Action & ObjectStatus & 7rem & -- & Action type badge (C=Create, U=Update, X=Delete) \\ \hline
  Environment & ObjectStatus & 5.5rem & -- & Target environment (DEV/QAS/PRD) \\ \hline
  Plant & Input/Text & 8rem & Yes & Plant identifier (value help available) \\ \hline
  Sending WH & Input/Text & 10rem & Yes & Sending warehouse code (value help) \\ \hline
  Receiving WH & Input/Text & 10rem & Yes & Receiving warehouse code (must differ from SendWh, BR-001) \\ \hline
  Transport Mode & Input/Text & 9rem & Yes & Transportation mode code (value help) \\ \hline
  Is Allowed & CheckBox & 6rem & -- & Whether this route is permitted (boolean) \\ \hline
  Inspector & Input/Text & 10rem & -- & Inspector user ID (composite value help: UserId OR Fullname) \\ \hline
  Version & Text & 4.5rem & -- & Version number (read-only) \\ \hline
  Change Note & Input/Text & flexible & -- & Description of change (max 255 chars) \\ \hline
\end{tabularx}
\caption{MM Routes --- UI Field Descriptions}
\label{tab:ui-mm-fields}
\end{table}

The MM Routes configuration module manages warehouse-to-warehouse material transfer routing rules within the Materials Management domain. In complex manufacturing and distribution environments, materials are frequently moved between warehouses within the same plant or across plants, and each route requires specific transportation modes, inspector assignments, and permission flags. A misconfigured route can result in materials being sent to an unauthorized warehouse, bypassing quality inspection checkpoints, or using an incorrect transportation mode that increases logistics costs or violates safety regulations.

The most distinctive validation rule in the MM Routes module is BR-001: the Sending Warehouse (SendWh) must differ from the Receiving Warehouse (ReceiveWh). This rule is enforced at both the frontend level (the \texttt{\_validateRows()} method checks for equality before allowing save) and the backend level (the \texttt{zcl\_gsp26\_rule\_validator} class performs the same check during the approve action). Additional validations include: PlantId, SendWh, ReceiveWh, and TransMode are all mandatory fields; the composite business key (PlantId + SendWh + ReceiveWh) must be unique within the request; and the Inspector field, while optional, must reference a valid user ID when provided. The Inspector value help dialog features a composite filter that searches across both UserId and Fullname fields simultaneously, allowing users to find inspectors by either their system ID or their display name.

The Excel import for MM Routes maps the following columns: ``Plant'' or ``PlantId'' or ``plant\_id'' maps to PlantId; ``Sending Warehouse'' or ``SendWh'' or ``send\_wh'' or ``From Warehouse'' maps to SendWh; ``Receiving Warehouse'' or ``ReceiveWh'' or ``receive\_wh'' or ``To Warehouse'' maps to ReceiveWh; ``Transport Mode'' or ``TransMode'' or ``trans\_mode'' maps to TransMode; ``Is Allowed'' or ``IsAllowed'' or ``Allowed'' maps to IsAllowed (parsed as boolean: ``X'', ``true'', ``yes'', ``1'' all evaluate to true); ``Inspector'' or ``InspectorId'' or ``inspector\_id'' maps to InspectorId; and ``Change Note'' or ``Note'' maps to ChangeNote. The boolean parsing for the IsAllowed field uses the \texttt{parseBool()} utility function in \texttt{ExcelImport.js}, which handles multiple common representations of boolean values in spreadsheets.

The draft choreography for MM Routes is identical to the FI Limit pattern, using the same POST-Activate sequence for new rows and Edit-PATCH-Activate sequence for modified rows. One module-specific consideration is the boolean IsAllowed field, which is rendered as a CheckBox control in edit mode and an ObjectStatus indicator in display mode. During the save process, the controller converts the CheckBox state to the SAP ABAP boolean representation (``X'' for true, empty string for false) before including it in the OData PATCH payload.

\paragraph{c. MM Safe Stock Configuration Screen}

\imgplaceholder{INSERT: Screenshot --- MM Safe Stock. Same layout. 7-column table: Action, Env, Plant, MatGroup, MinQty, Version, ChangeNote. Uses RAP Draft Choreography (POST then Activate per row). Value helps for Plant, MatGroup. MatGroup validated against master data. MinQty must be at least 0. Duplicate key check: PlantId+MatGroup.}

\begin{table}[H]
\centering\small
\begin{tabularx}{\textwidth}{|l|l|l|c|X|}
  \hline
  \rowcolor{tableheadpeach}
  \textbf{Field} & \textbf{Type} & \textbf{Width} & \textbf{Req.} & \textbf{Description} \\ \hline
  Action & ObjectStatus & 7rem & -- & Action type badge (C=Create, U=Update, X=Delete) \\ \hline
  Environment & ObjectStatus & 5.5rem & -- & Target environment (DEV/QAS/PRD) \\ \hline
  Plant & Input/Text & 8rem & Yes & Plant identifier (value help available) \\ \hline
  Material Group & Input/Text & 10rem & Yes & Material group code (value help, validated against master data) \\ \hline
  Min Quantity & Input (Number) & 9rem & Yes & Minimum safe stock quantity (must be >= 0) \\ \hline
  Version & Text & 4.5rem & -- & Version number (read-only) \\ \hline
  Change Note & Input/Text & flexible & -- & Description of change (max 255 chars) \\ \hline
\end{tabularx}
\caption{MM Safe Stock --- UI Field Descriptions}
\label{tab:ui-mmss-fields}
\end{table}

The MM Safe Stock configuration module defines minimum inventory levels for material groups at specific plants within the Materials Management domain. Safe stock (also known as safety stock) is a critical parameter in inventory management: it represents the minimum quantity of a material group that should be maintained in stock at a given plant to buffer against demand variability and supply chain disruptions. Setting safe stock levels too low risks stockouts and production delays; setting them too high ties up working capital in excess inventory and increases warehousing costs.

Validation rules for the MM Safe Stock module include: PlantId and MatGroup are mandatory fields; the MinQty value must be a non-negative number (BR-003), as negative safety stock has no physical meaning; the MatGroup value is validated against master data through the \texttt{ZI\_VH\_MATLGRP} value help CDS view, ensuring that only recognized material group codes can be used; and the composite business key (PlantId + MatGroup) must be unique within the request. The MM Safe Stock module is notable for using the RAP Draft Choreography pattern most explicitly among the four configuration apps, as its simpler field structure (only 3 editable business fields compared to 4--9 in other modules) makes the draft handling flow more visible in the code.

The Excel import for MM Safe Stock maps columns as follows: ``Plant'' or ``PlantId'' or ``plant\_id'' maps to PlantId; ``Material Group'' or ``MatGroup'' or ``mat\_group'' or ``Material Grp'' maps to MatGroup; ``Min Quantity'' or ``MinQty'' or ``min\_qty'' or ``Minimum Quantity'' or ``Safety Stock'' maps to MinQty; and ``Change Note'' or ``Note'' maps to ChangeNote. Because the MM Safe Stock module has fewer fields than other modules, the Excel import process is faster and less error-prone, making it particularly suitable for bulk data entry scenarios where inventory planners need to update safe stock levels for dozens or hundreds of plant-material group combinations simultaneously.

The draft choreography for MM Safe Stock follows the same pattern as the other modules but with a simpler payload structure. Each save operation involves fewer OData PATCH fields (PlantId, MatGroup, MinQty, ChangeNote compared to the 8--14 fields in FI or SD modules), resulting in smaller HTTP request payloads and faster save operations. The module also demonstrates the reusability of the \texttt{ConfMngBaseController} pattern: despite having a different number of fields, different validation rules, and different value help configurations than the other modules, the core lifecycle management (CSRF token handling, request header patching, submit flow, email notification) is entirely inherited from the base controller.

\paragraph{d. SD Price Configuration Screen}

\imgplaceholder{INSERT: Screenshot --- SD Price. Grid Table (sap.ui.table.Table). 14 columns: Action, Env, Branch, CustGrp, MatlGrp, MaxDiscount, Currency, MinOrderVal, ValidFrom, ValidTo, Approver, Version, ChangeNote. DatePicker for ValidFrom/ValidTo. 5 value help dialogs. Validation: ValidFrom must be on or before ValidTo, MaxDiscount 0--100 percent, MinOrderVal at least 0. Mandatory: BranchId, CustGroup, Currency. Duplicate key: BranchId+CustGroup+MaterialGrp.}

\begin{table}[H]
\centering\small
\begin{tabularx}{\textwidth}{|l|l|l|c|X|}
  \hline
  \rowcolor{tableheadpeach}
  \textbf{Field} & \textbf{Type} & \textbf{Width} & \textbf{Req.} & \textbf{Description} \\ \hline
  Action & ObjectStatus & 7rem & -- & Action type badge (C=Create, U=Update, X=Delete) \\ \hline
  Environment & ObjectStatus & 5.5rem & -- & Target environment (DEV/QAS/PRD) \\ \hline
  Branch & Input/Text & 8rem & Yes & Branch identifier (value help available) \\ \hline
  Customer Group & Input/Text & 10rem & Yes & Customer group code (value help) \\ \hline
  Material Group & Input/Text & 10rem & -- & Material group code (value help) \\ \hline
  Max Discount & Input (Number) & 8rem & -- & Maximum discount percentage (0--100, BR-012) \\ \hline
  Currency & Input/Text & 7rem & Yes & Currency code (VND/USD/EUR, value help) \\ \hline
  Min Order Value & Input (Number) & 9rem & -- & Minimum order value (must be >= 0) \\ \hline
  Valid From & DatePicker & 9rem & -- & Start date of validity period (ISO 8601) \\ \hline
  Valid To & DatePicker & 9rem & -- & End date of validity period (must be >= ValidFrom, BR-002) \\ \hline
  Approver & Input/Text & 9rem & -- & Approver group code (value help) \\ \hline
  Version & Text & 4.5rem & -- & Version number (read-only) \\ \hline
  Change Note & Input/Text & flexible & -- & Description of change (max 255 chars) \\ \hline
\end{tabularx}
\caption{SD Price --- UI Field Descriptions}
\label{tab:ui-sd-fields}
\end{table}


The SD Price configuration module manages sales pricing rules within the Sales and Distribution domain. Pricing rules in SD define the maximum discount percentages, minimum order values, validity periods, and approver assignments for specific combinations of branches, customer groups, and material groups. These parameters directly affect revenue recognition and customer contract enforcement: an incorrectly configured maximum discount could allow sales representatives to offer discounts beyond authorized levels, while an expired validity period could cause the system to apply outdated pricing terms to new orders.

The SD Price module has the most complex validation rule set among the four configuration apps. Key rules include: BranchId, CustGroup, and Currency are mandatory fields; the MaxDiscount value must be between 0 and 100 inclusive (BR-012), as a discount exceeding 100\% would imply paying the customer; the MinOrderVal must be a non-negative number (BR-003); the ValidFrom date must be on or before the ValidTo date (BR-002), ensuring that the validity period is logically consistent; and the composite business key (BranchId + CustGroup + MaterialGrp) must be unique within the request. The date validation is performed by comparing the DatePicker values using JavaScript Date object comparison, and the error message clearly identifies which row and which date field is in violation. The SD Price app uses a Grid Table (\texttt{sap.ui.table.Table}) instead of the responsive \texttt{sap.m.Table} used by the other three modules, because the 14-column layout requires the fixed-column rendering and horizontal scrolling capabilities that only the Grid Table provides.

The Excel import for SD Price maps the following columns: ``Branch'' or ``BranchId'' or ``branch\_id'' maps to BranchId; ``Customer Group'' or ``CustGroup'' or ``cust\_group'' or ``CustGrp'' maps to CustGroup; ``Material Group'' or ``MaterialGrp'' or ``material\_grp'' or ``MatlGrp'' maps to MaterialGrp; ``Max Discount'' or ``MaxDiscount'' or ``max\_discount'' or ``Discount'' maps to MaxDiscount; ``Currency'' or ``Curr'' maps to Currency; ``Min Order Value'' or ``MinOrderVal'' or ``min\_order\_val'' or ``MinOrder'' maps to MinOrderVal; ``Valid From'' or ``ValidFrom'' or ``valid\_from'' or ``Start Date'' maps to ValidFrom; ``Valid To'' or ``ValidTo'' or ``valid\_to'' or ``End Date'' maps to ValidTo; ``Approver'' or ``ApproverGrp'' or ``approver\_grp'' maps to ApproverGrp; and ``Change Note'' or ``Note'' maps to ChangeNote. Date fields in the Excel file are parsed using the SheetJS date utilities, which handle multiple date formats including ISO 8601 strings, Excel serial numbers, and common locale-specific formats (dd/MM/yyyy, MM/dd/yyyy).

The draft choreography for SD Price is the most complex among the four modules due to the larger number of fields and the presence of date fields that require special serialization. Date values from the DatePicker controls are converted to ISO 8601 format (yyyy-MM-dd) before being included in the OData PATCH payload, and the response dates from the backend are converted back to locale-specific display format for rendering in the table cells. The SD Price controller (approximately 1,208 lines) is the largest of the four module controllers, reflecting the additional complexity of managing 14 columns, 5 value help dialogs, and 2 date pickers within a single DynamicPage view.

\subsection{Manager App}

\paragraph{a. Manager Dashboard}

\imgplaceholder{INSERT: Screenshot --- Manager Dashboard. DynamicPage with gradient header. 4 KPI tiles (Pending=orange RadialMicroChart, Total=neutral NumericContent, Approved=green, Rejected=red). Animated counters via requestAnimationFrame (900ms ease-out). Skeleton shimmer loader during fetch. Two chart cards: Status donut chart (sap.viz), Module column chart. Recent Pending table (top 10 SUBMITTED, ordered by CreatedAt desc).}

\paragraph{b. Manager List Report}

\imgplaceholder{INSERT: Screenshot --- Manager List. Fiori Elements ListReport. Inline stats bar: 6 GenericTags (Total, Pending, Approved, Rejected, Approval percent, Rejection percent). Selection variants: All, Submitted, Approved, Rejected, Draft. Custom columns: ModuleColumn (InfoLabel), EnvColumn (InfoLabel), StatusColumn (InfoLabel with color scheme). Notification bell with 30s polling.}

\paragraph{c. Manager Object Page}

\imgplaceholder{INSERT: Screenshot --- Manager Object Page. Fiori Elements ObjectPage. Header: ReqTitle, CreatedBy, StatusDP. Actions: Approve (green), Reject (red). AuditTrailButton for cross-app nav. Items table with Compare button per row showing OldValue vs NewValue.}


\subsection{IT Admin App}

\paragraph{a. IT Admin My Inbox}

\imgplaceholder{INSERT: Screenshot --- IT Admin FCL layout. FlexibleColumnLayout (TwoColumnsMidExpanded). Begin column (33 percent): Sticky header with KPI badges (DEV/QAS/PRD counts), notification bell, EnvCompare button. Module TabBar (MM/MMSS/FI/SD). Filter strip: search, ENV segment, date range. Request list (single-select). Mid column (67 percent): Default = Audit Trail. On request select: Detail panel with IconTabBar (Preview/Journey/Info tabs). Footer: Promote, Rollback, Close buttons.}

\paragraph{b. Journey UI and Promote/Rollback}

\imgplaceholder{INSERT: Screenshot --- Journey Panel. 3 HTML cards (DEV, QAS, PRD) connected by lines. Each card: env label, status badge (Active/Promoted/Rolled Back), user, date. Current node: colored background + checkmark. Rollback button on each node. Inline rollback confirm panel (orange/red border, target env displayed, Confirm and Cancel buttons).}

\paragraph{c. Environment Compare Dialog}

\imgplaceholder{INSERT: Screenshot --- EnvCompare Dialog. Toolbar: Module dropdown, Compare button, Diffs Only toggle. Legend: green=Identical, orange=Different, red=Missing. Table: Record Key, Field, DEV value+status, QAS value+status, PRD value+status, Sync indicator. Version row per record. Growing table (200 threshold).}

\paragraph{d. Audit Diff Dialog}

\imgplaceholder{INSERT: Screenshot --- AuditDiff Dialog. Title: Diff followed by tableName. Context badges: action, module, env, table, user, time. Two-column grid: Left=Old values table (red highlight for changed), Right=New values table (green highlight for changed).}


% ------------------------------------------------------------
\section{Non-Functional Requirements}

\subsection{External Interfaces}

\begin{table}[H]
\centering\small
\begin{tabularx}{\textwidth}{|l|l|l|X|}
  \hline
  \rowcolor{tableheadpeach}
  \textbf{Interface} & \textbf{Protocol} & \textbf{Direction} & \textbf{Description} \\ \hline
  SAP OData V4 (8 svc) & HTTPS   & Bi-dir   & All CRUD, workflow actions, queries \\ \hline
  Email Service         & REST    & Outbound & POST to Render.com mail service \\ \hline
  SheetJS Library       & JS      & Local    & Excel file parsing (bundled) \\ \hline
  Browser localStorage  & JS API  & Local    & Notification state, nav mode \\ \hline
  SAP Ushell            & JS API  & Local    & User info retrieval \\ \hline
\end{tabularx}
\caption{External Interfaces}
\label{tab:external-interfaces}
\end{table}

\subsection{Quality Attributes}

\begin{table}[H]
\centering\small
\begin{tabularx}{\textwidth}{|l|X|X|}
  \hline
  \rowcolor{tableheadpeach}
  \textbf{Attribute} & \textbf{Requirement} & \textbf{Measurement} \\ \hline
  Performance    & Page load $<$ 3s; table render $<$ 1s for 100 rows    & DevTools Network \\ \hline
  Responsiveness & Notification poll $\leq$ 60s                          & Config constant \\ \hline
  Usability      & SAP Fiori Design Guidelines; $\leq$3 clicks           & Visual review \\ \hline
  Reliability    & Draft saves prevent data loss; works without email     & Manual test \\ \hline
  Security       & CSRF tokens; row-level DCL; role-based visibility      & Auth test \\ \hline
  Maintainability& Shared base controller for 4 apps; generic import     & Code review \\ \hline
\end{tabularx}
\caption{Quality Attributes}
\label{tab:quality-attributes}
\end{table}


% ------------------------------------------------------------
\section{Requirement Appendix}

\subsection{Business Rules}

\begin{longtable}{|l|p{3.5cm}|p{5.5cm}|l|}
  \hline
  \rowcolor{tableheadpeach}
  \textbf{Code} & \textbf{Rule} & \textbf{Description} & \textbf{Enforced} \\ \hline
  \endhead \hline \endfoot
  BR-001 & SendWh $\neq$ ReceiveWh       & Sending WH must differ from receiving WH (MM Routes) & FE + BE \\ \hline
  BR-002 & ValidFrom $\leq$ ValidTo       & Valid-from on or before valid-to (SD Price)           & FE \\ \hline
  BR-003 & Numeric $\geq$ 0               & AutoApprLim, MinOrderVal, MaxDiscount, MinQty         & FE + BE \\ \hline
  BR-004 & Reject reason required         & Manager must provide reason when rejecting             & FE + BE \\ \hline
  BR-005 & Promote DEV$\rightarrow$QAS$\rightarrow$PRD & Cannot skip environments during promotion & BE \\ \hline
  BR-006 & Rollback PRD$\rightarrow$QAS$\rightarrow$DEV & Must rollback higher envs first          & FE + BE \\ \hline
  BR-007 & Valid status transitions        & 9 transitions allowed (state machine)                 & BE \\ \hline
  BR-008 & Submit $\geq$1 item             & Cannot submit empty request                           & BE \\ \hline
  BR-009 & Approve validates items         & Required fields, no duplicates, source exists          & BE \\ \hline
  BR-010 & REJECTED is terminal            & No further actions on rejected requests                & BE \\ \hline
  BR-011 & Header immutable after approval & Cannot modify APPROVED/REJECTED headers                & BE \\ \hline
  BR-012 & MaxDiscount $\leq$ 100\%        & Maximum discount cannot exceed 100 percent             & FE \\ \hline
  BR-013 & Duplicate business key check    & No duplicate composite keys within same request        & FE \\ \hline
\end{longtable}

\subsection{Common Requirements}

\begin{itemize}[leftmargin=2em]
  \item All frontend apps support \texttt{sap-client} URL parameter (default: ``324'')
  \item Navigation mode determined by localStorage \texttt{conf-mng-nav-mode} or hostname auto-detection
  \item Notification persistence: localStorage, TTL 7 days for read items, max 50 notifications
  \item Email notifications: fire-and-forget; app unaffected by mail service failure
  \item All OData writes require CSRF token via \texttt{X-CSRF-Token: Fetch} header
  \item Date fields: ISO 8601 internally (\texttt{yyyy-MM-dd}), locale-specific display externally
\end{itemize}

\subsection{Application Messages List}

\begin{longtable}{|l|l|p{7cm}|l|}
  \hline
  \rowcolor{tableheadpeach}
  \textbf{Code} & \textbf{Type} & \textbf{Message} & \textbf{Context} \\ \hline
  \endhead \hline \endfoot
  MSG-001 & Error   & Expense Type is required                           & FI validation \\ \hline
  MSG-002 & Error   & GL Account is required                             & FI validation \\ \hline
  MSG-003 & Error   & Sending WH and Receiving WH must be different      & MM BR-001 \\ \hline
  MSG-004 & Error   & Valid From must be on or before Valid To            & SD BR-002 \\ \hline
  MSG-005 & Error   & Approval Limit must be 0 or greater                & FI BR-003 \\ \hline
  MSG-006 & Warning & Please provide a reason for this request            & Create request \\ \hline
  MSG-007 & Success & Request created. Edit what you need, then Save.     & After create \\ \hline
  MSG-008 & Success & X change(s) saved                                  & After save \\ \hline
  MSG-009 & Success & Request submitted for approval                      & After submit \\ \hline
  MSG-010 & Error   & Cannot fetch CSRF token                            & Network error \\ \hline
  MSG-011 & Info    & Changes are part of a change request                & Banner \\ \hline
  MSG-012 & Error   & Completed request cannot be changed                 & Edit approved \\ \hline
  MSG-013 & Warning & Rollback ENV first before rolling back ENV2         & BR-006 \\ \hline
  MSG-014 & Success & Promoted to ENV                                     & After promote \\ \hline
  MSG-015 & Success & Rolled back successfully                            & After rollback \\ \hline
  MSG-016 & Error   & Max Discount cannot exceed 100\%                    & SD BR-012 \\ \hline
  MSG-017 & Error   & Duplicate entry --- key already exists in request   & BR-013 \\ \hline
  MSG-018 & Error   & Material Group does not exist                       & MMSS validation \\ \hline
\end{longtable}

\subsection{Other Requirements}

\begin{itemize}[leftmargin=2em]
  \item \textbf{Browser Support:} Google Chrome (latest), Microsoft Edge (latest)
  \item \textbf{Responsive Design:} Desktop and tablet (SAP Fiori content densities: compact, cozy)
  \item \textbf{Internationalization:} English UI via i18n resource bundles; Vietnamese notifications in backend
  \item \textbf{Accessibility:} Standard SAP UI5 features (ARIA labels, keyboard navigation)
\end{itemize}


% ============================================================
\clearpage
\chapter{Software Design Description}

\section{System Design}

\subsection{System Architecture}

The system follows a strict four-layer architecture, separating concerns between presentation, service exposure, business logic, and data persistence. Figure~\ref{fig:arch} illustrates the overall system architecture.

\imgplaceholder{INSERT: System Architecture Diagram (4 layers) --- PRESENTATION LAYER: 8 SAP UI5 apps (Catalog :8080, Request :8081, FI :8084, MM :8082, MMSS :8085, SD :8083, Manager, ITAdmin). Shared modules: ConfMngBaseController, ExcelImport, NotificationService. --- SERVICE LAYER: SAP Gateway with 8 Service Definitions (ZSD\_*), 8 Service Bindings (ZUI\_*), Consumption Views (ZC\_*). Protocol: OData V4. --- BUSINESS LOGIC LAYER: Behavior Handler zbp\_ir\_conf\_req\_h (8 actions). 6 Rule Classes: rule\_status, rule\_validator, rule\_snapshot, rule\_version, rule\_writer, conf\_api. 6 DCL files for row-level security. --- DATA LAYER: 51 CDS Interface Views (ZI\_*), 31 Database Tables. SAP HANA. --- EXTERNAL: Mail Service (Node.js on Render.com), fire-and-forget POST /api/notify to Gmail API.}

\begin{enumerate}[leftmargin=2em]
  \item \textbf{Presentation Layer:} Eight SAP UI5 Fiori applications, each communicating with the backend exclusively through OData V4 services. Shared patterns are implemented via \texttt{ConfMngBaseController} (URL parameter parsing, CSRF token management, request lifecycle, email notification), \texttt{ExcelImport.js} (config-driven Excel parsing), and \texttt{NotificationService.js} (polling with localStorage persistence).

  \item \textbf{Service Layer:} SAP Gateway exposes 8 OData V4 service bindings. Each service definition projects interface CDS views into consumption views with appropriate UI annotations. Table~\ref{tab:odata-services} lists all service bindings.

  \item \textbf{Business Logic Layer:} The core behavior handler \texttt{zbp\_ir\_conf\_req\_h} (approximately 1,810 lines) implements all 8 workflow actions. Six rule classes encapsulate reusable business logic: state machine validation, field validation, snapshot/rollback, version numbering, audit logging, and a cached read API for integration scenarios.

  \item \textbf{Data Layer:} 31 database tables organized into five groups (Request, Config Request, Config Main, Master/Reference, Value Help), accessed through 51 CDS interface views with row-level security via 6 DCL files.
\end{enumerate}

\subsubsection{OData V4 Draft Handling}

The OData V4 draft handling mechanism is a central architectural element of the system, enabling users to save work-in-progress configuration changes without immediately committing them to the active dataset. In the SAP RAP framework, draft handling is declared in the behavior definition using the \texttt{with draft;} annotation, which automatically generates a draft table (suffixed with \texttt{\_D}) for each entity. The draft protocol operates on the concept of \texttt{IsActiveEntity}: when this flag is \texttt{true}, the entity represents the committed (active) version; when \texttt{false}, it represents the draft (work-in-progress) version.

The configuration apps in this system use an explicit draft choreography rather than relying on the SAP Fiori Elements automatic draft management. This design decision was made because the apps need to overlay request-specific changes on top of the main configuration data --- a pattern that the standard Fiori Elements draft handling does not support. The explicit choreography involves three key operations: (1) \textbf{Create Draft:} A POST request to the entity set creates a new draft instance with \texttt{IsActiveEntity=false}. The response includes the server-generated keys (ItemId, ReqItemId). (2) \textbf{Activate:} A POST request to the \texttt{draftActivate} action converts the draft into an active entity, running all backend validations defined in the behavior definition's \texttt{determination} and \texttt{validation} blocks. (3) \textbf{Edit:} A POST request to the \texttt{draftEdit} action on an existing active entity creates a draft copy that can be modified via PATCH requests before being activated.

\subsubsection{CSRF Token Flow}

All write operations (POST, PATCH, DELETE) to the SAP OData V4 services require a valid Cross-Site Request Forgery (CSRF) token. The frontend applications implement a centralized token management strategy in the \texttt{ConfMngBaseController}. When a write operation is needed, the controller first sends a GET request to the OData service metadata endpoint with the header \texttt{X-CSRF-Token: Fetch}. The server responds with the token in the \texttt{X-CSRF-Token} response header, which is then cached in the controller instance variable \texttt{\_sCsrfToken}. All subsequent POST and PATCH requests include this cached token in their \texttt{X-CSRF-Token} request header. If a token expires during a long editing session (SAP tokens typically have a 30-minute lifetime), the server returns a 403 Forbidden response, and the controller automatically fetches a new token and retries the failed request. This retry logic is implemented in the \texttt{fetchCsrfToken()} method of the base controller and is shared across all eight frontend applications.

\subsubsection{Cross-App Navigation}

The system comprises eight independent SAP UI5 applications, each running on its own localhost port during development (Catalog on 8080, Request on 8081, MM Routes on 8082, SD Price on 8083, FI Limit on 8084, MM Safe Stock on 8085) or deployed as separate BSP applications on the SAP server. Cross-app navigation is implemented using a combination of URL parameters and localStorage state management. When one app needs to navigate to another (e.g., the Catalog app navigating to the FI Limit config app after creating a request), it constructs the target URL with query parameters including \texttt{ReqId}, \texttt{ConfId}, \texttt{ModuleId}, \texttt{ConfName}, \texttt{TargetCds}, and \texttt{sap-client}. The navigation mode --- whether to use localhost ports or SAP BSP paths --- is determined by a localStorage flag (\texttt{conf-mng-nav-mode}) that can be set to ``local'' or ``bsp''. If this flag is not set, the system auto-detects the mode by checking whether the current hostname is ``localhost'' or a SAP server hostname. This dual-mode navigation design allows developers to test the full cross-app workflow on their local machines while ensuring seamless operation when deployed to the SAP server.

\subsubsection{Fire-and-Forget Email Notification Pattern}

The email notification architecture follows the fire-and-forget pattern, which was chosen for three specific reasons. First, the external mail service is hosted on Render.com's free tier, which has a cold-start latency of up to 30 seconds when the service has not received a request recently. Awaiting the response would create an unacceptable delay in the user interface. Second, email delivery is inherently unreliable --- network issues, Gmail API rate limits, or temporary service outages can cause delivery failures at any time. By not depending on successful email delivery, the core application workflow remains robust and predictable regardless of external service availability. Third, the fire-and-forget pattern simplifies error handling in the frontend: the \texttt{sendMailNotification()} method in \texttt{ConfMngBaseController} wraps the \texttt{fetch()} call in a try-catch block that silently swallows all errors, including network timeouts, HTTP 4xx/5xx responses, and JSON parsing failures. The application continues to function identically whether the mail service is online, offline, or experiencing errors.

\begin{table}[H]
\centering\small
\begin{tabularx}{\textwidth}{|l|l|X|}
  \hline
  \rowcolor{tableheadpeach}
  \textbf{Service Binding} & \textbf{Root Entity} & \textbf{Purpose} \\ \hline
  zui\_conf\_req       & ZC\_CONF\_REQ\_H    & Request CRUD and workflow actions \\ \hline
  zui\_conf\_catalog   & ZC\_CONF\_CATALOG   & Catalog browser \\ \hline
  zui\_fi\_limit\_conf  & ZC\_FI\_LIMIT\_CONF  & FI configurations (draft-enabled) \\ \hline
  zui\_mm\_route\_conf  & ZC\_MM\_ROUTE\_CONF  & MM routes (draft-enabled) \\ \hline
  zui\_mm\_safe\_stock  & ZC\_MM\_SAFE\_STOCK  & MMSS configurations (draft-enabled) \\ \hline
  zsd\_sd\_price\_conf  & ZC\_SD\_PRICE\_CONF  & SD prices (draft-enabled) \\ \hline
  zui\_audit\_log      & ZC\_AUDIT\_LOG      & Audit trail (read-only) \\ \hline
  zui\_conf\_mng\_app   & ---                 & Config Management App \\ \hline
\end{tabularx}
\caption{OData V4 Service Bindings}
\label{tab:odata-services}
\end{table}


\subsection{Package Diagram}

\paragraph{Backend Package Structure}

\imgplaceholder{INSERT: Backend Package Diagram --- ZGSP26SAP06 (ABAP Package): src/core/ (44 files: CDS Interface Views ZI\_* 20, Behavior Definitions *.bdef 7, Behavior Implementations ZBP\_* 6, Abstract Entities ZABSTR\_*, Value Help Views ZI\_VH\_* 11); src/srv/ (35 files: Consumption Views ZC\_* 18, Service Definitions ZSD\_* 8, Main Views ZC\_*\_MAIN 4); src/rule/ (6 files: rule\_status, rule\_validator, rule\_snapshot, rule\_version, rule\_writer, conf\_api); src/sec/ (6 DCL files); src/req/ (1 notification provider); src/testdata/ (20 files: 3 test classes, 11 seed programs, 5 clean programs).}

The backend package \texttt{ZGSP26SAP06} is organized into six sub-directories, each with a clear responsibility:

\begin{itemize}[leftmargin=2em]
  \item \textbf{src/core/} (44 files): Contains all CDS interface views (\texttt{ZI\_*}), behavior definitions (\texttt{*.bdef}), behavior handler implementations (\texttt{ZBP\_*}), abstract entities for action parameters (\texttt{ZABSTR\_*}), and value help views (\texttt{ZI\_VH\_*}).
  \item \textbf{src/srv/} (35 files): Contains consumption/projection views (\texttt{ZC\_*}), service definitions (\texttt{ZSD\_*}), and read-only main views (\texttt{ZC\_*\_MAIN}) used by frontend apps to load current configuration data.
  \item \textbf{src/rule/} (6 files): Six reusable rule classes that encapsulate business logic independently from the RAP behavior handler.
  \item \textbf{src/sec/} (6 files): Data Control Language (DCL) files implementing row-level security. The request header DCL restricts access based on \texttt{CreatedBy = \$session.user} for Key Users, while Managers and IT Admins have broader access.
  \item \textbf{src/req/} (1 file): SAP Gateway notification provider (\texttt{zcl\_gsp26\_req\_notif\_prv}) that sends push notifications using Vietnamese templates.
  \item \textbf{src/testdata/} (20 files): Three test classes for automated testing, 11 seed programs for consistent test data, and 5 cleanup programs for test isolation.
\end{itemize}

\paragraph{Frontend Package Structure}

\imgplaceholder{INSERT: Frontend Package Diagram (typical config app) --- webapp/: Component.js (bootstrap), manifest.json (routes, models, dataSources), index.html; ext/main/: Main.controller.js (app-specific CRUD), Main.view.xml (DynamicPage layout), Main.css; ext/shared/: ConfMngBaseController.js (shared base), ExcelImport.js (generic import), NotificationService.js (polling), NotificationPopover.fragment.xml; lib/xlsx.min.js (SheetJS bundled); annotations/, i18n/, localService/.}

Each frontend application follows a consistent directory structure. The shared modules (\texttt{ConfMngBaseController}, \texttt{ExcelImport}, \texttt{NotificationService}) are duplicated across the four configuration apps (FI, MM, MMSS, SD) to maintain independence while sharing the same implementation patterns.


\section{Database Design}

The system uses 31 database tables organized into five groups. The entity relationship is centered on the Request Header (\texttt{ZCONFREQH}) which links to Request Items (\texttt{ZCONFREQI}) and module-specific configuration request tables via \texttt{ReqId}.

\imgplaceholder{INSERT: Entity Relationship Diagram --- Full ERD showing relationships between all 31 tables across 5 groups: Request tables (ZCONFREQH, ZCONFREQI with drafts), Config Request tables (ZFILIMITREQ, ZMMROUTECONF\_REQ, ZMMSAFESTOCK\_REQ, ZSD\_PRICE\_REQ linked by ReqId), Config Main tables (ZFILIMITCONF, ZMMROUTECONF, ZMMSAFESTOCK, ZSD\_PRICE\_CONF linked by SourceItemId), Master tables (ZCONFCATALOG, ZCONFFIELDDEF, ZAUDITLOG, ZUSERROLE, ZENVDEF), and 11 Value Help tables.}

\subsection{Request Tables}

\begin{table}[H]
\centering\small
\begin{tabularx}{\textwidth}{|l|l|X|}
  \hline
  \rowcolor{tableheadpeach}
  \textbf{Table} & \textbf{Key} & \textbf{Purpose} \\ \hline
  ZCONFREQH    & ReqId (UUID) + EnvId & Request headers: status, title, reason, approval info \\ \hline
  ZCONFREQI    & ReqItemId (UUID)     & Request items: action type, target env, version \\ \hline
  ZCONFREQH\_D & (draft table)        & Draft table for request headers \\ \hline
  ZCONFREQI\_D & (draft table)        & Draft table for request items \\ \hline
\end{tabularx}
\caption{Request Tables}
\label{tab:request-tables}
\end{table}

\noindent\textbf{ZCONFREQH --- Request Header Fields:}

\begin{longtable}{|l|l|c|p{6.5cm}|}
  \hline
  \rowcolor{tableheadpeach}
  \textbf{Field} & \textbf{Type} & \textbf{Key} & \textbf{Description} \\ \hline
  \endhead \hline \endfoot
  req\_id         & UUID (RAW16)  & PK & Request identifier \\ \hline
  env\_id         & CHAR10        & PK & Environment (DEV/QAS/PRD) \\ \hline
  conf\_id        & UUID          &    & Link to configuration catalog \\ \hline
  module\_id      & CHAR10        &    & Module identifier (FI, MM, SD) \\ \hline
  req\_title      & CHAR100       &    & Request display title \\ \hline
  description    & CHAR255       &    & Detailed description \\ \hline
  reason         & CHAR255       &    & Reason for change \\ \hline
  status         & CHAR20        &    & DRAFT, SUBMITTED, APPROVED, REJECTED, ACTIVE, ROLLED\_BACK \\ \hline
  reject\_reason  & CHAR255       &    & Manager's rejection reason \\ \hline
  created\_by     & CHAR12        &    & SAP user ID of creator \\ \hline
  created\_at     & TIMESTAMPL    &    & Creation timestamp \\ \hline
  approved\_by    & CHAR12        &    & Approving manager \\ \hline
  approved\_at    & TIMESTAMPL    &    & Approval timestamp \\ \hline
  rejected\_by    & CHAR12        &    & Rejecting manager \\ \hline
  rejected\_at    & TIMESTAMPL    &    & Rejection timestamp \\ \hline
  changed\_by     & CHAR12        &    & Last modifier \\ \hline
  changed\_at     & TIMESTAMPL    &    & Last modification timestamp \\ \hline
\end{longtable}

\subsection{Configuration Request Tables}

Each module has a request table storing the delta (changes proposed in a request). All four tables share a common structure with module-specific business fields.

\begin{table}[H]
\centering\small
\begin{tabularx}{\textwidth}{|l|l|X|}
  \hline
  \rowcolor{tableheadpeach}
  \textbf{Table} & \textbf{Module} & \textbf{Business Key Fields} \\ \hline
  ZFILIMITREQ        & FI   & ExpenseType, GlAccount \\ \hline
  ZMMROUTECONF\_REQ  & MM   & PlantId, SendWh, ReceiveWh \\ \hline
  ZMMSAFESTOCK\_REQ  & MMSS & PlantId, MatGroup \\ \hline
  ZSD\_PRICE\_REQ    & SD   & BranchId, CustGroup, MaterialGrp \\ \hline
\end{tabularx}
\caption{Module-Specific Configuration Request Tables}
\label{tab:config-req-tables}
\end{table}

\noindent\textbf{Common fields across all configuration request tables:}

\begin{table}[H]
\centering\small
\begin{tabularx}{\textwidth}{|l|l|X|}
  \hline
  \rowcolor{tableheadpeach}
  \textbf{Field} & \textbf{Type} & \textbf{Description} \\ \hline
  req\_id          & UUID   & Link to request header \\ \hline
  req\_item\_id    & UUID   & Link to request item \\ \hline
  item\_id         & UUID   & Unique line identifier \\ \hline
  action\_type     & CHAR1  & C = Create, U = Update, X = Delete \\ \hline
  source\_item\_id & UUID   & Reference to original config line (for U/X) \\ \hline
  env\_id          & CHAR10 & Target environment \\ \hline
  version\_no      & INT4   & Version number \\ \hline
  old\_*           & varies & Previous values for each business field \\ \hline
  change\_note     & CHAR255& Change description \\ \hline
\end{tabularx}
\caption{Common Fields in Configuration Request Tables}
\label{tab:common-req-fields}
\end{table}

\noindent\textbf{ZFILIMITREQ --- FI Limit-Specific Fields:}

\begin{table}[H]
\centering\small
\begin{tabularx}{\textwidth}{|l|l|X|}
  \hline
  \rowcolor{tableheadpeach}
  \textbf{Field} & \textbf{Type} & \textbf{Description} \\ \hline
  expense\_type & CHAR20 & Expense category code \\ \hline
  gl\_account & CHAR10 & General Ledger account number \\ \hline
  auto\_appr\_lim & DECIMAL(15,2) & Automatic approval limit amount \\ \hline
  currency & CHAR5 & Currency code (VND, USD, EUR) \\ \hline
  old\_expense\_type & CHAR20 & Previous expense type (for U/X actions) \\ \hline
  old\_gl\_account & CHAR10 & Previous GL account \\ \hline
  old\_auto\_appr\_lim & DECIMAL(15,2) & Previous approval limit \\ \hline
  old\_currency & CHAR5 & Previous currency \\ \hline
\end{tabularx}
\caption{ZFILIMITREQ --- FI Limit Fields}
\label{tab:fi-fields}
\end{table}

\noindent\textbf{ZMMROUTECONF\_REQ --- MM Routes-Specific Fields:}

\begin{table}[H]
\centering\small
\begin{tabularx}{\textwidth}{|l|l|X|}
  \hline
  \rowcolor{tableheadpeach}
  \textbf{Field} & \textbf{Type} & \textbf{Description} \\ \hline
  plant\_id & CHAR10 & Plant identifier \\ \hline
  send\_wh & CHAR10 & Sending warehouse code \\ \hline
  receive\_wh & CHAR10 & Receiving warehouse code (must differ from send\_wh) \\ \hline
  trans\_mode & CHAR10 & Transportation mode code \\ \hline
  is\_allowed & CHAR1 & Whether route is permitted (X = true, empty = false) \\ \hline
  inspector\_id & CHAR12 & Inspector user ID \\ \hline
  old\_plant\_id & CHAR10 & Previous plant (for U/X actions) \\ \hline
  old\_send\_wh & CHAR10 & Previous sending warehouse \\ \hline
  old\_receive\_wh & CHAR10 & Previous receiving warehouse \\ \hline
  old\_trans\_mode & CHAR10 & Previous transportation mode \\ \hline
  old\_is\_allowed & CHAR1 & Previous allowed flag \\ \hline
  old\_inspector\_id & CHAR12 & Previous inspector \\ \hline
\end{tabularx}
\caption{ZMMROUTECONF\_REQ --- MM Routes Fields}
\label{tab:mm-fields}
\end{table}

\noindent\textbf{ZMMSAFESTOCK\_REQ --- MM Safe Stock-Specific Fields:}

\begin{table}[H]
\centering\small
\begin{tabularx}{\textwidth}{|l|l|X|}
  \hline
  \rowcolor{tableheadpeach}
  \textbf{Field} & \textbf{Type} & \textbf{Description} \\ \hline
  plant\_id & CHAR10 & Plant identifier \\ \hline
  mat\_group & CHAR20 & Material group code (validated against master data) \\ \hline
  min\_qty & DECIMAL(13,3) & Minimum safe stock quantity (must be >= 0) \\ \hline
  old\_plant\_id & CHAR10 & Previous plant (for U/X actions) \\ \hline
  old\_mat\_group & CHAR20 & Previous material group \\ \hline
  old\_min\_qty & DECIMAL(13,3) & Previous minimum quantity \\ \hline
\end{tabularx}
\caption{ZMMSAFESTOCK\_REQ --- MM Safe Stock Fields}
\label{tab:mmss-fields}
\end{table}

\noindent\textbf{ZSD\_PRICE\_REQ --- SD Price-Specific Fields:}

\begin{table}[H]
\centering\small
\begin{tabularx}{\textwidth}{|l|l|X|}
  \hline
  \rowcolor{tableheadpeach}
  \textbf{Field} & \textbf{Type} & \textbf{Description} \\ \hline
  branch\_id & CHAR10 & Branch identifier \\ \hline
  cust\_group & CHAR20 & Customer group code \\ \hline
  material\_grp & CHAR20 & Material group code \\ \hline
  max\_discount & DECIMAL(5,2) & Maximum discount percentage (0--100) \\ \hline
  min\_order\_val & DECIMAL(15,2) & Minimum order value \\ \hline
  approver\_grp & CHAR20 & Approver group code \\ \hline
  currency & CHAR5 & Currency code (VND, USD, EUR) \\ \hline
  valid\_from & DATS & Validity start date \\ \hline
  valid\_to & DATS & Validity end date (must be >= valid\_from) \\ \hline
  old\_branch\_id & CHAR10 & Previous branch (for U/X actions) \\ \hline
  old\_cust\_group & CHAR20 & Previous customer group \\ \hline
  old\_material\_grp & CHAR20 & Previous material group \\ \hline
  old\_max\_discount & DECIMAL(5,2) & Previous maximum discount \\ \hline
  old\_min\_order\_val & DECIMAL(15,2) & Previous minimum order value \\ \hline
  old\_approver\_grp & CHAR20 & Previous approver group \\ \hline
  old\_currency & CHAR5 & Previous currency \\ \hline
  old\_valid\_from & DATS & Previous validity start date \\ \hline
  old\_valid\_to & DATS & Previous validity end date \\ \hline
\end{tabularx}
\caption{ZSD\_PRICE\_REQ --- SD Price Fields}
\label{tab:sd-fields}
\end{table}

\subsection{Configuration Main Tables and Master Tables}

\begin{table}[H]
\centering\small
\begin{tabularx}{\textwidth}{|l|l|X|}
  \hline
  \rowcolor{tableheadpeach}
  \textbf{Table} & \textbf{Module/Type} & \textbf{Purpose} \\ \hline
  ZFILIMITCONF   & FI (Main)     & Approved FI limit configurations by environment \\ \hline
  ZMMROUTECONF   & MM (Main)     & Approved MM route configurations by environment \\ \hline
  ZMMSAFESTOCK   & MMSS (Main)   & Approved safe stock configurations by environment \\ \hline
  ZSD\_PRICE\_CONF & SD (Main)   & Approved SD price configurations by environment \\ \hline
  ZCONFCATALOG   & Master        & Configuration catalog definitions \\ \hline
  ZCONFFIELDDEF  & Master        & Field definitions per catalog (name, label, required, type) \\ \hline
  ZAUDITLOG      & Audit         & Immutable audit trail (action, old/new JSON, timestamp) \\ \hline
  ZUSERROLE      & Security      & User role assignments (KEY USER, MANAGER, IT ADMIN) \\ \hline
  ZENVDEF        & Reference     & Environment definitions (DEV, QAS, PRD) \\ \hline
\end{tabularx}
\caption{Configuration Main Tables and Master Tables}
\label{tab:main-master-tables}
\end{table}

\noindent\textbf{Value Help Tables (11):} Branch, Currency, CustGroup, ExpenseType, GlAccount, Inspector, MatlGrp, Module, TransMode, Warehouse, PlantUnit. These tables provide dropdown/select dialog data for frontend value help dialogs.


\section{Detailed Design}

\subsection{Request Workflow}

\subsubsection{State Machine}

The request lifecycle is governed by a finite state machine with 7 states and 9 valid transitions, implemented in \texttt{zcl\_gsp26\_rule\_status}. Figure~\ref{fig:state-machine} illustrates the complete state diagram.

\imgplaceholder{INSERT: State Machine Diagram --- 7 states: DRAFT, SUBMITTED, APPROVED, REJECTED (terminal), ACTIVE, PROMOTED, ROLLED\_BACK. Transitions: DRAFT to SUBMITTED (submit), SUBMITTED to APPROVED (approve), SUBMITTED to REJECTED (reject), APPROVED to ACTIVE (apply), APPROVED to PROMOTED (promote), APPROVED to ROLLED\_BACK (rollback), ACTIVE to PROMOTED (promote), ACTIVE to ROLLED\_BACK (rollback), PROMOTED to ROLLED\_BACK (rollback).}

\begin{table}[H]
\centering\small
\begin{tabularx}{\textwidth}{|c|l|l|l|X|}
  \hline
  \rowcolor{tableheadpeach}
  \textbf{No} & \textbf{From} & \textbf{To} & \textbf{Action} & \textbf{Triggered By} \\ \hline
  1 & DRAFT & SUBMITTED & submit & Key User \\ \hline
  2 & SUBMITTED & APPROVED & approve & Manager \\ \hline
  3 & SUBMITTED & REJECTED & reject & Manager \\ \hline
  4 & APPROVED & ACTIVE & apply & IT Admin \\ \hline
  5 & APPROVED & PROMOTED & promote & IT Admin \\ \hline
  6 & APPROVED & ROLLED\_BACK & rollback & IT Admin \\ \hline
  7 & ACTIVE & PROMOTED & promote & IT Admin \\ \hline
  8 & ACTIVE & ROLLED\_BACK & rollback & IT Admin \\ \hline
  9 & PROMOTED & ROLLED\_BACK & rollback & IT Admin \\ \hline
\end{tabularx}
\caption{Valid State Transitions}
\label{tab:state-transitions}
\end{table}

\noindent\textbf{State Descriptions:}

\begin{itemize}[leftmargin=2em]
  \item \textbf{DRAFT:} The initial state when a request is first created. In this state, the Key User can add, edit, and delete configuration lines, import data from Excel, and save drafts. The request is only visible to its creator. No approval actions are available. The request can be submitted once it contains at least one configuration item (BR-008).

  \item \textbf{SUBMITTED:} The request has been submitted for managerial review. In this state, the request becomes visible to all Managers in the Manager App. The Key User can no longer edit the request. The Manager can either approve or reject the request. This state represents the ``pending approval'' phase of the workflow.

  \item \textbf{APPROVED:} A Manager has reviewed and approved the request. The approval action triggers the creation of an approval snapshot (serialized JSON of the current configuration state) stored in \texttt{ZAUDITLOG}. This snapshot serves as the restore point for potential future rollbacks. The request is now available for IT Admin actions (apply, promote, rollback).

  \item \textbf{REJECTED:} A Manager has rejected the request with a mandatory reason. This is a terminal state --- no further workflow actions are available on this request. The rejection reason is stored in the \texttt{reject\_reason} field and displayed to the Key User. The Key User may choose to create a new request addressing the feedback, or resubmit the rejected request after editing.

  \item \textbf{ACTIVE:} The IT Administrator has applied the approved configuration changes to the DEV environment. The configuration data has been written to the module-specific main tables (ZFILIMITCONF, ZMMROUTECONF, ZMMSAFESTOCK, ZSD\_PRICE\_CONF) for the DEV environment. From this state, the request can be promoted to QAS or rolled back.

  \item \textbf{PROMOTED:} The configuration has been promoted to a higher environment (QAS or PRD). The promotion action copies the configuration data from the current environment to the next one in the pipeline. When promoted to PRD, the global version number is incremented. The request can be further promoted (if not yet at PRD) or rolled back.

  \item \textbf{ROLLED\_BACK:} The IT Administrator has reverted the configuration changes in one or more environments. The rollback action restores the previous configuration state from the approval snapshot stored in the audit log. Rollback order is strictly enforced: PRD must be rolled back before QAS, and QAS before DEV (BR-006). This ensures data consistency across environments.
\end{itemize}

The state machine is implemented as a simple lookup table in \texttt{zcl\_gsp26\_rule\_status}. The method \texttt{is\_transition\_valid\_by\_status()} accepts the current status and the target status as parameters and returns a boolean. This design avoids database access during validation, providing fast transition checks that can be called from any workflow action.

Any transition not listed in Table~\ref{tab:state-transitions} is considered invalid and will be rejected by the backend with an appropriate error message. This strict enforcement ensures that the request lifecycle follows the intended governance workflow at all times.

\subsubsection{Approve Flow --- Backend Validation Details}

The approve action is the most validation-intensive operation in the entire system. When a Manager clicks ``Approve'' in the Manager App, the backend behavior handler \texttt{zbp\_ir\_conf\_req\_h} executes a multi-step validation pipeline before committing the status change. Understanding this pipeline is essential for appreciating the data integrity guarantees provided by the system.

The first validation step is the status transition check: the \texttt{zcl\_gsp26\_rule\_status} class verifies that the current status is SUBMITTED and the target status is APPROVED. If the request is in any other status (e.g., DRAFT, REJECTED, or already APPROVED), the transition is rejected immediately. The second step is the item count validation: the handler queries \texttt{ZCONFREQI} to count the number of items linked to the request and rejects the approval if the count is zero (BR-008). The third step is the per-module field validation: for each configuration item, the \texttt{zcl\_gsp26\_rule\_validator} class reads the field definitions from \texttt{ZCONFFIELDDEF} for the relevant module and checks that all fields marked as \texttt{IsRequired = true} have non-empty values. For the FI module, this means ExpenseType, GlAccount, AutoApprLim, and Currency must all be populated. For the MM module, PlantId, SendWh, ReceiveWh, and TransMode are required. For MMSS, PlantId, MatGroup, and MinQty are mandatory. For SD, BranchId, CustGroup, and Currency are required.

The fourth validation step is duplicate business key detection: the validator checks that no two items within the same request share the same composite business key. For FI, the composite key is (ExpenseType + GlAccount); for MM, it is (PlantId + SendWh + ReceiveWh); for MMSS, it is (PlantId + MatGroup); and for SD, it is (BranchId + CustGroup + MaterialGrp). If duplicates are detected, the approval is rejected with an error message identifying the duplicate key values. The fifth step applies only to Update (ActionType=``U'') and Delete (ActionType=``X'') items: the validator verifies that the \texttt{SourceItemId} references an existing record in the corresponding main configuration table for the target environment. This ensures that the request is not attempting to update or delete a record that has already been removed by another request.

\subsubsection{Snapshot Mechanism}

The snapshot mechanism is the foundation of the system's rollback capability. When a request is approved, the \texttt{zcl\_gsp26\_rule\_snapshot} class creates a point-in-time snapshot of the configuration data that will be affected by the request. This snapshot captures the ``before'' state of the data, enabling precise restoration during rollback.

The snapshot creation process works as follows: for each configuration item in the request, the rule class reads the current values from the module-specific main table (e.g., \texttt{ZFILIMITCONF} for FI, \texttt{ZMMROUTECONF} for MM) for the target environment. These values are serialized into JSON format using the SAP standard \texttt{/ui2/cl\_json=>serialize()} method, which converts ABAP data structures into JSON strings with proper type handling for decimals, dates, and boolean fields. The serialized ``old data'' JSON is stored in the \texttt{old\_data} field of a new \texttt{ZAUDITLOG} entry, while the ``new data'' JSON (representing the proposed changes from the request items) is stored in the \texttt{new\_data} field. The audit log entry also records the action type (APPROVE), the module identifier, the environment, the requesting user, the approving user, and a precise timestamp.

During rollback, the \texttt{restore\_from\_snapshot()} method retrieves the most recent audit log entry for the given request and environment combination, deserializes the \texttt{old\_data} JSON back into an ABAP data structure using \texttt{/ui2/cl\_json=>deserialize()}, and writes the restored values back to the main configuration table. For items that were originally created (ActionType=``C''), the rollback deletes them from the main table. For items that were originally updated (ActionType=``U''), the rollback overwrites the current values with the snapshot values. For items that were originally deleted (ActionType=``X''), the rollback re-inserts them into the main table with their pre-deletion values. A new audit log entry with action type ROLLBACK is created, with the old and new data fields reversed compared to the original entry, providing a complete and reversible audit chain.

\subsubsection{Version Numbering Strategy}

Version numbering in the system serves two purposes: tracking the history of configuration changes and ensuring data consistency during promotion to the Production environment. The version number is a simple integer that is incremented only when configuration data is promoted to PRD --- changes applied to DEV and promoted to QAS retain the version number assigned during creation.

The \texttt{zcl\_gsp26\_rule\_version} class implements the \texttt{get\_next\_version()} method, which determines the next available version number for a given module and environment combination. This method uses dynamic SQL (ABAP \texttt{SELECT MAX(version\_no) FROM (table\_name)}) to query the maximum existing version number from the target main table, allowing a single method to work across all four module tables without hardcoding table names. The dynamic table name is derived from the module identifier using a simple mapping: FI maps to \texttt{ZFILIMITCONF}, MM maps to \texttt{ZMMROUTECONF}, MMSS maps to \texttt{ZMMSAFESTOCK}, and SD maps to \texttt{ZSD\_PRICE\_CONF}. The method adds 1 to the maximum version found and returns the result. If no records exist in the target table (i.e., this is the first promotion to PRD for this module), the method returns version 1.

This strategy ensures that PRD version numbers provide a meaningful chronological ordering of production-level configuration changes, while DEV and QAS versions remain at their initial values to avoid confusion with the production version sequence. Frontend applications display the version number in a read-only column, providing users with immediate visibility into which version of a configuration record they are viewing.

\subsubsection{Class Diagram}

\imgplaceholder{INSERT: Class Diagram --- Request Workflow. Main class: zbp\_ir\_conf\_req\_h (behavior handler) with methods: +approve(), +reject(reason), +submit(), +promote(), +apply(), +rollback(), +createRequest(params), +updateReason(params), +get\_instance\_authorizations(), +get\_instance\_features(), -set\_default\_and\_admin(), -validate\_before\_save(). Uses three rule classes: zcl\_gsp26\_rule\_status (+is\_transition\_valid\_by\_status(current, next) returns bool), zcl\_gsp26\_rule\_validator (+check\_catalog\_existence(), +check\_required\_fields(), +check\_field\_ranges(), +validate\_request\_item()), zcl\_gsp26\_rule\_snapshot (+create\_approve\_snapshot(), +restore\_from\_snapshot(), +increment\_version()). rule\_snapshot uses zcl\_gsp26\_rule\_writer (+log\_audit\_entry() serializes to JSON via /ui2/cl\_json).}

The behavior handler \texttt{zbp\_ir\_conf\_req\_h} serves as the central orchestrator for all workflow actions. It delegates specific business logic concerns to specialized rule classes, following the Single Responsibility Principle:

\begin{itemize}[leftmargin=2em]
  \item \textbf{zcl\_gsp26\_rule\_status:} Validates state transitions using a predefined state machine with 9 valid transitions. The method \texttt{is\_transition\_valid\_by\_status()} accepts current and target status as parameters, avoiding database access for better performance.
  \item \textbf{zcl\_gsp26\_rule\_validator:} Performs data validation including catalog existence checks, required field validation (driven by \texttt{ZCONFFIELDDEF}), field range checks, and composite request item validation.
  \item \textbf{zcl\_gsp26\_rule\_snapshot:} Manages configuration snapshots for the approve and promote workflows. The \texttt{create\_approve\_snapshot()} method serializes current configuration state to JSON and stores it in \texttt{ZAUDITLOG}. The \texttt{restore\_from\_snapshot()} method deserializes the JSON and restores previous values during rollback.
  \item \textbf{zcl\_gsp26\_rule\_version:} Provides dynamic version numbering via \texttt{get\_next\_version()}, which executes a \texttt{SELECT MAX(version\_no)} on the target table using dynamic SQL.
  \item \textbf{zcl\_gsp26\_rule\_writer:} Writes immutable audit log entries to \texttt{ZAUDITLOG}, serializing old and new data structures to JSON using \texttt{/ui2/cl\_json=>serialize()}.
  \item \textbf{zcl\_gsp26\_conf\_api:} Provides a cached read API for integration scenarios, with per-environment static caches for FI, MM, MMSS, and SD configurations.
\end{itemize}


\subsubsection{Sequence Diagram --- Submit and Approve Flow}

\imgplaceholder{INSERT: Sequence Diagram --- Submit: KeyUser clicks Submit, FE validates ReqId/Status=DRAFT, FE POST submit action to OData, zbp calls rule\_status.is\_valid(DRAFT,SUBMITTED)=true, zbp checks item count at least 1, zbp MODIFY Status=SUBMITTED, return. FE POST /api/notify event:SUBMITTED (fire-and-forget). FE shows success dialog. --- Approve: Manager clicks Approve, ManagerApp POST approve to OData, zbp calls rule\_status.is\_valid(SUBMITTED,APPROVED), zbp calls rule\_validator.validate\_request\_item() per item, zbp per-module validation (required fields, duplicates, source), zbp calls rule\_snapshot.create\_approve\_snapshot(req\_id) INSERT into ZAUDITLOG, zbp calls rule\_writer.log\_audit\_entry(APPROVE), zbp MODIFY Status=APPROVED/ApprovedBy/ApprovedAt, return. FE POST /api/notify event:APPROVED.}


\subsubsection{Sequence Diagram --- Promote Flow}

\imgplaceholder{INSERT: Sequence Diagram --- ITAdmin clicks Promote, FE determines nextEnv (DEV to QAS or QAS to PRD), confirmation dialog. FE POST promote to OData. zbp validates status APPROVED/ACTIVE, calculates nextEnv. LOOP per module (MM, MMSS, FI, SD): SELECT from source env, UPDATE/INSERT into target env. IF target=PRD: rule\_version.get\_next\_version(). rule\_writer.log\_audit\_entry(PROMOTE). END LOOP. zbp MODIFY EnvId=nextEnv, return. FE POST /api/notify event:PROMOTED.}


\subsection{Configuration CRUD}

\subsubsection{Class Diagram}

\imgplaceholder{INSERT: Class Diagram --- Frontend Config Apps. Abstract class ConfMngBaseController with abstract methods: getModuleId(), getConfName(), getTargetCds(). Public methods: onEdit(), onNavBack(), onCreateRequest(), onSubmitRequest(). Protected methods: getUrlParams(), fetchCsrfToken(), fetchFieldDef(), fetchRequestHeader(), fetchReqItem(), patchRequestHeader(), sendMailNotification(), extractODataError(), dbToPascal(), statusToState(). --- Extended by 4 concrete controllers: FI Main (1058 lines), MM Main (1157 lines), MMSS Main (991 lines, RAP Draft Choreography), SD Main (1208 lines, 14 cols + date validation). --- Utility classes: ExcelImport.js (+parseWorkbook, -mapHeaders, -transformRow, -parseBool, -normalize), NotificationService.js (+init, +markRead, +markAllRead, +clearAll, +destroy, -poll).}

The frontend architecture follows the Template Method pattern, where \texttt{ConfMngBaseController} defines the skeleton of the request lifecycle (create, submit, CSRF management, email notification) and each module-specific controller overrides abstract methods and implements its own CRUD logic, validation rules, and value help dialogs.


\subsubsection{Sequence Diagram --- Save Flow (Draft Choreography)}

\imgplaceholder{INSERT: Sequence Diagram --- User clicks Save Draft. MainController calls validateRows(allRows). If errors: MessageBox.error, return. MainController GET CSRF token from OData. LOOP per changed row: IF state=new: POST /ModuleConf (create draft, IsActiveEntity=false), POST Activate (draft becomes active), store reqItemId. IF state=modified with existing reqItemId: POST Edit (creates draft copy), PATCH draft with changed fields, POST Activate. END LOOP. MainController POST updateReason (patch title+reason). Set editMode=false. Reload table with overlay. Toast: X change(s) saved.}

This ``Draft Choreography'' pattern is a key architectural decision. Rather than using the SAP Fiori Elements draft handling (which manages drafts automatically via the framework), the configuration apps manage drafts explicitly through direct OData calls. This provides fine-grained control over the save process, particularly for the overlay loading pattern where main table data is merged with request-specific changes.


\subsection{Notification System}

\subsubsection{Sequence Diagram --- Polling}

\imgplaceholder{INSERT: Sequence Diagram --- ON INIT: NotifSvc loads seen/list/lastPoll from localStorage. NotifSvc GET ZC\_CONF\_REQ\_H top 100 (initial catchup). Compare statuses vs mSeen map. Update JSONModel (notifications[], unreadCount). Badge shows count. --- EVERY 15--60 SECONDS: Fetch top 100 requests. Per request: compare Status vs mSeen[ReqId]. IF changed (not DRAFT to SUBMITTED): buildNotif(). Prepend to notifications, increment unread. Prune expired (read items older than 7 days). Cap at 50. Save to localStorage.}


\subsection{Email Service}

\subsubsection{Sequence Diagram --- Email Notification}

\imgplaceholder{INSERT: Sequence Diagram --- FE POST /api/notify to MailService (Render.com). Headers: x-api-key, Content-Type application/json. Body: event, reqId, reqTitle, module, envId, triggeredBy. MailService validates API key, builds HTML email template, sends via nodemailer to Gmail API. Recipient receives email. NOTE: FE does not await response (fire-and-forget). All errors silently caught, app continues normally.}

The email notification service is intentionally decoupled from the core application. The fire-and-forget pattern ensures that email delivery failures (including Render.com cold-start delays of up to 30 seconds) never block or degrade the user experience. The mail service is a stateless Node.js application that accepts notification payloads via REST API and forwards them to Gmail using the \texttt{nodemailer} library.


% ============================================================
%       CHAPTER V.  SOFTWARE TESTING DOCUMENTATION
% ============================================================
\clearpage
\chapter{Software Testing Documentation}

\section{Scope of Testing}

\textbf{In Scope:}
\begin{itemize}[leftmargin=2em]
  \item Functional testing of all 18 use cases
  \item Unit testing of Excel import modules (4 apps $\times$ QUnit)
  \item Unit testing of backend rule classes (ABAP Unit)
  \item Integration testing of basic navigation flows (OPA5)
  \item System testing of end-to-end workflows (manual + ABAP test classes)
  \item Cross-app navigation testing
  \item Role-based access control testing
  \item Data integrity testing across 3 environments
\end{itemize}

\textbf{Out of Scope:}
\begin{itemize}[leftmargin=2em]
  \item Performance/load testing (limited by shared server)
  \item Automated UI regression testing
  \item Security penetration testing
  \item Mobile device testing
\end{itemize}

\begin{table}[H]
\centering\small
\begin{tabularx}{\textwidth}{|l|l|X|X|}
  \hline
  \rowcolor{tableheadpeach}
  \textbf{Level} & \textbf{In-Charge} & \textbf{Focus} & \textbf{Acceptance Criteria} \\ \hline
  Unit        & Developers     & Individual modules (ExcelImport, Rule classes) & All cases pass; edge cases covered \\ \hline
  Integration & Developers     & App navigation, OData service calls             & Apps launch; routes work \\ \hline
  System      & QA / All       & End-to-end workflows across all apps            & All 18 UCs executable \\ \hline
  UAT         & Supervisor     & Real-world scenario validation                  & Supervisor confirms \\ \hline
\end{tabularx}
\caption{Testing Levels}
\label{tab:testing-levels}
\end{table}


The testing approach was tailored to the specific characteristics of each application layer. For the backend rule classes (\texttt{zcl\_gsp26\_rule\_status}, \texttt{zcl\_gsp26\_rule\_validator}, \texttt{zcl\_gsp26\_rule\_snapshot}), ABAP Unit tests were written to validate both positive and negative scenarios. The status machine tests, for example, covered all 9 valid transitions to confirm they return \texttt{true} and a representative set of invalid transitions (such as DRAFT to APPROVED, REJECTED to ACTIVE, and ROLLED\_BACK to SUBMITTED) to confirm they return \texttt{false}. For the frontend ExcelImport modules, QUnit tests were written to validate header mapping with various alias combinations, number parsing with different locale formats (comma vs. period as decimal separator), boolean parsing for the MM Routes IsAllowed field, and workbook parsing with both valid and malformed data.

Several test scenarios proved particularly challenging during the project. Testing the rollback workflow across all three environments required careful orchestration: a request had to be created, approved, applied to DEV, promoted to QAS, and promoted to PRD before the rollback sequence (PRD first, then QAS, then DEV) could be tested. Any test failure in the middle of this sequence left the test data in an inconsistent state, requiring manual cleanup before the test could be re-executed. The team addressed this by creating a dedicated ABAP test class (\texttt{zcl\_test\_end\_to\_end}) that automated the entire lifecycle and included a cleanup method in the \texttt{teardown()} section to restore the database to its pre-test state regardless of where a failure occurred.

Testing concurrent draft editing --- where two Key Users attempt to edit the same configuration simultaneously --- revealed that the OData V4 draft protocol provides natural concurrency protection through its \texttt{IsActiveEntity} mechanism: when one user enters edit mode (creating a draft copy), a second user attempting to edit the same entity receives an error because the draft lock is already held. This behavior was verified through manual testing with two browser sessions logged in as different Key Users.

Testing Excel import with malformed data required a comprehensive set of test files covering edge cases such as: empty spreadsheets (0 data rows), spreadsheets with headers but no data rows, spreadsheets with unrecognized column headers, spreadsheets with extra columns not defined in the alias dictionary, spreadsheets with numeric values stored as text strings, spreadsheets with date values in multiple formats, and spreadsheets with more than 500 rows (triggering the confirmation dialog). These test files were stored in the project's test data directory and referenced in the QUnit test suite for automated regression testing.

Test data management was a critical concern given the shared SAP server environment. The team developed 11 seed programs (\texttt{zseed\_*.prog.abap}) that populate the database tables with consistent, known test data before each testing session. Each seed program uses the ABAP \texttt{MODIFY} statement with internal tables to insert or update records in a specific table, ensuring idempotent execution (running the seed program multiple times produces the same result). Complementing the seed programs, 5 cleanup programs (\texttt{zclean\_*.prog.abap}) remove test-generated data (such as draft requests, test audit log entries, and temporary configuration records) without affecting the seed data. The typical testing workflow was: (1) run cleanup programs to remove residual test data, (2) run seed programs to establish a known baseline, (3) execute test cases, (4) record results, and (5) run cleanup programs again to restore the baseline for the next tester.


\section{Test Strategy}

\subsection{Testing Types}

\begin{table}[H]
\centering\small
\begin{tabularx}{\textwidth}{|l|X|X|X|}
  \hline
  \rowcolor{tableheadpeach}
  \textbf{Type} & \textbf{Objective} & \textbf{Technique} & \textbf{Completion Criteria} \\ \hline
  Functional   & Verify each UC works as specified           & Black-box, BVA              & All main + key alt flows pass \\ \hline
  Regression   & Ensure fixes don't break existing features  & Re-run affected test cases  & No previously passing tests fail \\ \hline
  Usability    & Verify Fiori UX compliance                  & Heuristic evaluation        & No critical usability issues \\ \hline
  Security     & Verify authorization and data protection    & Role-based access, CSRF     & Unauthorized access blocked \\ \hline
  Integration  & Verify cross-app communication              & E2E navigation, OData chains& All cross-app flows complete \\ \hline
\end{tabularx}
\caption{Testing Types}
\label{tab:testing-types}
\end{table}

\subsection{Test Levels}

\begin{itemize}[leftmargin=2em]
  \item \textbf{Unit Level:} QUnit for frontend ExcelImport (\texttt{parseNumber}, \texttt{mapHeaders}, \texttt{transformRow}, \texttt{parseWorkbook}); ABAP Unit for \texttt{zcl\_test\_rule\_classes} (status transitions, field validation, snapshot/rollback)
  \item \textbf{Integration Level:} OPA5 journeys for app launch and basic navigation; \texttt{zcl\_test\_rap\_flow} for RAP action integration
  \item \textbf{System Level:} Manual E2E testing with real SAP data; \texttt{zcl\_test\_end\_to\_end} for automated full workflow (create $\rightarrow$ approve $\rightarrow$ promote)
\end{itemize}

\subsection{Supporting Tools}

\begin{table}[H]
\centering\small
\begin{tabularx}{\textwidth}{|l|X|l|}
  \hline
  \rowcolor{tableheadpeach}
  \textbf{Tool} & \textbf{Purpose} & \textbf{Used At} \\ \hline
  QUnit              & JavaScript unit testing      & FE ExcelImport \\ \hline
  OPA5               & UI5 integration testing      & FE navigation \\ \hline
  ABAP Unit          & ABAP class testing           & BE rule classes \\ \hline
  ADT Test Runner    & Execute ABAP tests           & Eclipse IDE \\ \hline
  Chrome DevTools    & Debug, network inspection    & FE development \\ \hline
  SAP UI5 Diagnostics& Performance tracing          & FE optimization \\ \hline
  Postman            & OData service testing        & API verification \\ \hline
\end{tabularx}
\caption{Test Supporting Tools}
\label{tab:test-tools}
\end{table}


\section{Test Plan}

\subsection{Human Resources}

\begin{table}[H]
\centering\small
\begin{tabularx}{\textwidth}{|l|X|l|}
  \hline
  \rowcolor{tableheadpeach}
  \textbf{Role} & \textbf{Responsibility} & \textbf{Member} \\ \hline
  Test Lead        & Plan and coordinate testing activities       & \placeholder{Name} \\ \hline
  Unit Tester (FE) & Write and execute QUnit / OPA5 tests        & \placeholder{Name} \\ \hline
  Unit Tester (BE) & Write and execute ABAP Unit tests           & \placeholder{Name} \\ \hline
  System Tester    & Execute E2E workflows manually              & \placeholder{Name} \\ \hline
  UAT Coordinator  & Coordinate with supervisor for acceptance   & \placeholder{Name} \\ \hline
\end{tabularx}
\caption{Test Human Resources}
\label{tab:test-hr}
\end{table}

\subsection{Test Environment}

\begin{table}[H]
\centering\small
\begin{tabularx}{\textwidth}{|l|X|}
  \hline
  \rowcolor{tableheadpeach}
  \textbf{Component} & \textbf{Details} \\ \hline
  SAP Server  & s40lp1.ucc.cit.tum.de, Client 324 \\ \hline
  Local Dev   & localhost ports 8080--8085 (ui5 serve) \\ \hline
  Browser     & Google Chrome / Microsoft Edge (latest) \\ \hline
  Test Users  & 3 users with KEY USER, MANAGER, IT ADMIN roles \\ \hline
  Seed Data   & 11 seed programs (\texttt{zseed\_*.prog}) for consistent test data \\ \hline
  Cleanup     & 5 clean programs (\texttt{zclean\_*.prog}) for test isolation \\ \hline
\end{tabularx}
\caption{Test Environment}
\label{tab:test-env}
\end{table}

\subsection{Test Milestones}

\begin{table}[H]
\centering\small
\begin{tabularx}{\textwidth}{|X|l|l|}
  \hline
  \rowcolor{tableheadpeach}
  \textbf{Milestone} & \textbf{Target Date} & \textbf{Status} \\ \hline
  Unit test cases written               & \placeholder{Date} & \placeholder{Status} \\ \hline
  Unit tests executed                   & \placeholder{Date} & \placeholder{Status} \\ \hline
  Integration tests completed           & \placeholder{Date} & \placeholder{Status} \\ \hline
  System testing round 1                & \placeholder{Date} & \placeholder{Status} \\ \hline
  Defect fixes verified                 & \placeholder{Date} & \placeholder{Status} \\ \hline
  System testing round 2 (regression)   & \placeholder{Date} & \placeholder{Status} \\ \hline
  UAT completed                         & \placeholder{Date} & \placeholder{Status} \\ \hline
\end{tabularx}
\caption{Test Milestones}
\label{tab:test-milestones}
\end{table}


\section{Test Cases}

Test cases are documented in separate Excel files, organized by testing level:

\begin{itemize}[leftmargin=2em]
  \item \textbf{Unit Test Cases:} \texttt{Report5\_Unit Test.xlsx}
    \begin{itemize}
      \item Frontend: \texttt{ExcelImport.qunit.js} --- 15+ test cases covering \texttt{parseNumber}, \texttt{mapHeaders}, \texttt{transformRow}, and \texttt{parseWorkbook} with both valid inputs and edge cases (empty data, invalid headers, numeric parsing).
      \item Backend: \texttt{zcl\_test\_rule\_classes} --- tests for status transitions (valid and invalid), field validation against \texttt{ZCONFFIELDDEF}, and snapshot creation/restoration.
    \end{itemize}

  \item \textbf{Integration / System / Acceptance Test Cases:} \texttt{Report5\_Test Report.xlsx}
    \begin{itemize}
      \item Organized by UC code (UC-001 through UC-018)
      \item Each test case includes: Test ID, UC reference, precondition, test steps, expected result, actual result, pass/fail status
      \item Cross-app navigation scenarios tested separately
    \end{itemize}
\end{itemize}


\subsection{Sample Test Cases}

The following tables present representative test cases for key workflows. Complete test case documentation is provided in the Excel files referenced above.

\begin{table}[H]
\centering\small
\begin{tabularx}{\textwidth}{|l|l|X|X|l|}
  \hline
  \rowcolor{tableheadpeach}
  \textbf{TC-ID} & \textbf{UC} & \textbf{Steps} & \textbf{Expected Result} & \textbf{Status} \\ \hline
  TC-001 & UC-001 & 1. Open Catalog Detail for FI Limit. 2. Click ``Maintain via Request''. 3. Verify redirect to FI config app. & New request created with Status=DRAFT. FI config app opens in edit mode with current DEV data loaded. & \placeholder{P/F} \\ \hline
  TC-002 & UC-002 & 1. Click ``Add Line'' in FI config app. 2. Fill ExpenseType=TRAVEL, GlAccount=600100, Limit=5000, Currency=VND. 3. Click ``Save Draft''. & Row saved successfully. Toast: ``1 change(s) saved''. Row appears with green highlight (ActionType=C). & \placeholder{P/F} \\ \hline
  TC-003 & UC-003 & 1. Select an existing row in the table. 2. Click ``Delete''. 3. Confirm deletion dialog. 4. Click ``Save Draft''. & Row marked with red highlight (ActionType=X). After save, row persisted with deletion marker. & \placeholder{P/F} \\ \hline
  TC-004 & UC-004 & 1. Click ``Import from Excel''. 2. Select a valid .xlsx file with 5 rows and correct headers. 3. Verify imported rows. & Toast: ``5 rows imported (0 skipped)''. All 5 rows appear in table with green highlight and ActionType=C. & \placeholder{P/F} \\ \hline
  TC-005 & UC-004 & 1. Click ``Import from Excel''. 2. Select a .xlsx file with unrecognized column headers. & Error message: ``No recognizable headers found''. No rows imported. Table unchanged. & \placeholder{P/F} \\ \hline
  TC-006 & UC-006 & 1. With a DRAFT request containing 3 items, click ``Submit Request''. 2. Confirm in dialog. & Status changes to SUBMITTED. Success dialog shown. Email notification sent to managers. Navigation back to request list. & \placeholder{P/F} \\ \hline
  TC-007 & UC-007 & 1. Open a SUBMITTED request in Manager App. 2. Review items using Compare button. 3. Click ``Approve''. & Status changes to APPROVED. Approval snapshot created in ZAUDITLOG. Email sent to IT Admin and creator. & \placeholder{P/F} \\ \hline
  TC-008 & UC-008 & 1. Open a SUBMITTED request in Manager App. 2. Click ``Reject''. 3. Enter reason: ``Incorrect GL account''. 4. Confirm. & Status changes to REJECTED. Rejection reason saved. Email sent to request creator. Request becomes terminal. & \placeholder{P/F} \\ \hline
  TC-009 & UC-010 & 1. Open an APPROVED request in IT Admin App. 2. Click ``Promote''. 3. Confirm promotion to QAS. & Config data copied from DEV to QAS. Audit log entry created (action=PROMOTE). Email notification sent. & \placeholder{P/F} \\ \hline
  TC-010 & UC-011 & 1. Open a request promoted to QAS. 2. Click ``Rollback'' on the QAS journey card. 3. Confirm rollback. & QAS config restored from snapshot. Status = ROLLED\_BACK. Audit log entry (action=ROLLBACK) with reversed old/new. & \placeholder{P/F} \\ \hline
  TC-011 & UC-011 & 1. Attempt to rollback DEV while QAS is still active (not rolled back). & Error: ``Rollback QAS first before rolling back DEV''. Rollback blocked (BR-006). & \placeholder{P/F} \\ \hline
  TC-012 & UC-014 & 1. Open Env Compare dialog in IT Admin App. 2. Select module ``FI''. 3. Click ``Compare''. 4. Toggle ``Diffs Only''. & Comparison table shows DEV/QAS/PRD values with color-coded sync status. Diffs Only toggle filters to differences only. & \placeholder{P/F} \\ \hline
\end{tabularx}
\caption{Sample Functional Test Cases}
\label{tab:sample-test-cases}
\end{table}

\noindent\textbf{Test Case Coverage Summary:}

The 12 sample test cases above cover the following critical workflows:
\begin{itemize}[leftmargin=2em]
  \item \textbf{Configuration CRUD:} TC-001 through TC-005 validate the create, edit, delete, and Excel import capabilities, including both positive and negative scenarios.
  \item \textbf{Approval Workflow:} TC-006 through TC-008 validate the submit, approve, and reject workflows, including the mandatory rejection reason and email notifications.
  \item \textbf{Environment Promotion:} TC-009 through TC-011 validate the promote and rollback workflows, including the strict rollback order enforcement (BR-006).
  \item \textbf{Supporting Features:} TC-012 validates the environment comparison feature used by IT Administrators to detect configuration drift.
\end{itemize}

The complete test suite contains \placeholder{N} test cases covering all 18 use cases, including boundary value analysis, error handling, cross-app navigation, and role-based access control scenarios. Detailed results are documented in the Excel test report files.


\section{Test Reports}

\begin{table}[H]
\centering\small
\begin{tabularx}{\textwidth}{|X|c|c|c|c|}
  \hline
  \rowcolor{tableheadpeach}
  \textbf{Test Type} & \textbf{Total} & \textbf{Pass} & \textbf{Fail} & \textbf{Pass Rate} \\ \hline
  Unit Test (FE --- QUnit)      & \placeholder{N} & \placeholder{N} & \placeholder{N} & \placeholder{\%} \\ \hline
  Unit Test (BE --- ABAP Unit)  & \placeholder{N} & \placeholder{N} & \placeholder{N} & \placeholder{\%} \\ \hline
  Integration Test (OPA5)       & \placeholder{N} & \placeholder{N} & \placeholder{N} & \placeholder{\%} \\ \hline
  System Test (E2E)             & \placeholder{N} & \placeholder{N} & \placeholder{N} & \placeholder{\%} \\ \hline
  UAT                           & \placeholder{N} & \placeholder{N} & \placeholder{N} & \placeholder{\%} \\ \hline
  \textbf{Total}                & \placeholder{N} & \placeholder{N} & \placeholder{N} & \placeholder{\%} \\ \hline
\end{tabularx}
\caption{Test Execution Summary}
\label{tab:test-results}
\end{table}

\imgplaceholder{INSERT: Test Execution Summary Chart --- Bar chart showing Pass (green) and Fail (red) counts per test type. X-axis: Unit(FE), Unit(BE), Integration, System, UAT. Y-axis: Count of test cases.}


% ============================================================
\clearpage
\chapter{Release Package \& User Guides}

\section{Deliverable Package}

Table~\ref{tab:deliverable-package} lists all deliverable items included in the final release of the System Configuration Management project.

\begin{table}[H]
\centering\small
\begin{tabularx}{\textwidth}{|c|l|X|c|}
  \hline
  \rowcolor{tableheadpeach}
  \textbf{No} & \textbf{Item} & \textbf{Description} & \textbf{Ver.} \\ \hline
  1 & Frontend Source Code  & 8 SAP UI5 repos (GitHub: CONFIG-MNG-APP)         & 1.0 \\ \hline
  2 & Backend Source Code   & ABAP package ZGSP26SAP06 (126 files, abapGit)    & 1.0 \\ \hline
  3 & Mail Service Code     & Node.js notification service (Render.com)        & 1.0 \\ \hline
  4 & Final Report          & Capstone Project Document (this file)            & 1.0 \\ \hline
  5 & Test Cases            & Unit + System test Excel files                   & 1.0 \\ \hline
  6 & Presentation Slides   & Final defense presentation                       & 1.0 \\ \hline
  7 & Database Seed Scripts & 11 seed programs (\texttt{zseed\_*.prog.abap})   & 1.0 \\ \hline
  8 & Deployment Guide      & SAP BSP + Render.com deployment instructions     & 1.0 \\ \hline
  9 & Schedule / Tracking   & Sprint backlogs and tracking                     & Final \\ \hline
\end{tabularx}
\caption{Deliverable Package}
\label{tab:deliverable-package}
\end{table}


\section{Installation Guides}

\subsection{System Requirements}

Table~\ref{tab:sys-requirements} outlines the minimum system requirements for installing and running the application.

\begin{table}[H]
\centering\small
\begin{tabularx}{\textwidth}{|l|X|}
  \hline
  \rowcolor{tableheadpeach}
  \textbf{Component} & \textbf{Requirement} \\ \hline
  SAP Server           & SAP S/4HANA with RAP support, ABAP Platform $\geq$ 2020 \\ \hline
  ABAP Development     & Eclipse IDE with ADT plugin, abapGit installed \\ \hline
  Frontend Development & Node.js $\geq$ 18, npm $\geq$ 9, @ui5/cli \\ \hline
  Mail Service         & Node.js $\geq$ 18, Gmail account with App Password \\ \hline
  Browser              & Google Chrome or Microsoft Edge (latest) \\ \hline
  Network              & Access to SAP server (VPN if required) \\ \hline
\end{tabularx}
\caption{System Requirements}
\label{tab:sys-requirements}
\end{table}

\subsection{Installation Instructions}

\subsubsection{Backend Installation}

The backend ABAP code is managed via abapGit and deployed to the SAP S/4HANA server.

\begin{enumerate}[leftmargin=2em]
  \item \textbf{Clone abapGit repository:}
    \begin{itemize}
      \item Open abapGit in SAP GUI (transaction \texttt{ZABAPGIT}) or Eclipse ADT
      \item Create a new online repository pointing to:
      \item[] \texttt{https://github.com/CONFIG-MNG-APP/configuration-mgn-rap.git}
      \item Pull all objects into package \texttt{ZGSP26SAP06}
    \end{itemize}

  \item \textbf{Activate objects in dependency order:}
    \begin{itemize}
      \item Activate data elements and domains first
      \item Activate database table definitions
      \item Activate CDS interface views (\texttt{ZI\_*})
      \item Activate behavior definitions (\texttt{*.bdef}) and implementations (\texttt{ZBP\_*})
      \item Activate consumption/projection views (\texttt{ZC\_*})
      \item Activate service definitions (\texttt{ZSD\_*}) and service bindings (\texttt{ZUI\_*})
      \item Register OData V4 services in transaction \texttt{/IWFND/V4\_ADMIN}
    \end{itemize}

  \item \textbf{Seed test data:}
    \begin{itemize}
      \item Execute seed programs in the following order:
      \item[] \texttt{zseed\_confcatalog} $\rightarrow$ \texttt{zseed\_conffielddef} $\rightarrow$ \texttt{zseed\_userrole}
      \item[] $\rightarrow$ \texttt{zseed\_fi\_limit} $\rightarrow$ \texttt{zseed\_mm\_route\_conf} $\rightarrow$ \texttt{zseed\_mm\_safestock}
      \item[] $\rightarrow$ \texttt{zseed\_sd\_price\_conf} $\rightarrow$ value help seeds (\texttt{zseed\_fi\_vh}, etc.)
    \end{itemize}
\end{enumerate}

\subsubsection{Frontend Installation (per app)}

Each of the 8 frontend applications follows the same installation process.

\begin{enumerate}[leftmargin=2em]
  \item \textbf{Clone repository:}
\begin{lstlisting}
git clone https://github.com/CONFIG-MNG-APP/<repo-name>.git
cd <repo-name>
\end{lstlisting}

  \item \textbf{Install dependencies:}
\begin{lstlisting}
npm install
\end{lstlisting}

  \item \textbf{Local development:}
\begin{lstlisting}
npx ui5 serve -o index.html
\end{lstlisting}
    The application opens on its configured port:
    \begin{itemize}
      \item Catalog: 8080, Request: 8081, MM Routes: 8082
      \item SD Price: 8083, FI Limit: 8084, MM SafeStock: 8085
    \end{itemize}

  \item \textbf{SAP deployment (BSP):}
\begin{lstlisting}
npx ui5 build --all
\end{lstlisting}
    Deploy the \texttt{dist/} folder to the SAP BSP repository via Eclipse ADT or the ABAP file store endpoint \texttt{/sap/bc/adt/filestore/ui5-bsp/}.
\end{enumerate}

\subsubsection{Mail Service Installation}

The email notification service is a standalone Node.js application deployed on Render.com.

\begin{enumerate}[leftmargin=2em]
  \item \textbf{Clone and install:}
\begin{lstlisting}
git clone https://github.com/CONFIG-MNG-APP/mail-service.git
cd mail-service
npm install
\end{lstlisting}

  \item \textbf{Configure environment variables:}
    Create a \texttt{.env} file in the project root with the following variables:
\begin{lstlisting}
GMAIL_USER=your-email@gmail.com
GMAIL_APP_PASSWORD=your-app-password
API_KEY=abap19-mail-secret-key
PORT=3000
\end{lstlisting}

  \item \textbf{Local testing:}
\begin{lstlisting}
npm start
\end{lstlisting}
    The service starts on port 3000 and listens for POST requests at \texttt{/api/notify}.

  \item \textbf{Deploy to Render.com:}
    \begin{itemize}
      \item Connect the GitHub repository to a new Render Web Service
      \item Set the environment variables in the Render dashboard
      \item Auto-deploy triggers on each push to the main branch
      \item Production URL: \texttt{https://abap19-mail-206.onrender.com}
    \end{itemize}
\end{enumerate}


\section{User Manual}

\subsection{Overview}

The System Configuration Management application provides a complete workflow for managing configuration changes across SAP modules. The workflow involves three user roles working in sequence, as illustrated in Figure~\ref{fig:workflow-overview}.

\imgplaceholder{INSERT: Application Workflow Overview Diagram --- Three columns for KEY USER, MANAGER, and IT ADMIN. Flow: 1. Browse Catalog (Key User), 2. Create Request (Key User), 3. Edit Config Lines (Key User), 4. Submit Request (Key User, arrow to Manager), 5. Review and Approve (Manager, arrow to IT Admin), 6. Apply to DEV (IT Admin), 7. Promote DEV to QAS to PRD (IT Admin). Optional: Rollback at any promoted environment.}

The following subsections provide step-by-step guides for the three primary workflows.


\subsection{Workflow 1: Create and Submit Configuration Request}

\textbf{Purpose:} Key Users create a change request to propose modifications to configuration values in any of the 4 modules (FI Limit, MM Routes, MM Safe Stock, SD Price).

\paragraph{Step 1: Open Configuration Catalog}
Navigate to the Catalog app (port 8080 in local development). The main page displays all available configurations with KPI cards showing counts by module (MM, SD, FI) and status (Total, Active, Inactive).

\imgplaceholder{INSERT: Screenshot --- Catalog List page. Shows DynamicPage with pinned header containing Module Info cards (MM, SD, FI with counts) and KPI cards (Total blue, Active green, Inactive grey). Content area has SegmentedButton filter (ALL/MM/SD/FI), SearchField, and a table with columns: Module, Config Name, Type, Description, Target CDS, Status. Growing table with 20 rows per batch.}

\paragraph{Step 2: Select Configuration and Click ``Maintain''}
Click on a configuration row to view its detail page. The detail page shows the configuration metadata, field definitions table, and three action buttons:
\begin{itemize}[leftmargin=2em]
  \item \textbf{View Current Config:} Opens the configuration data in read-only mode
  \item \textbf{View Requests:} Shows all existing requests for this configuration
  \item \textbf{Maintain via Request:} Creates a new change request and navigates to the configuration editor
\end{itemize}

Click the \textbf{``Maintain via Request''} button (emphasized blue button) to create a new change request.

\imgplaceholder{INSERT: Screenshot --- Catalog Detail page. Shows configuration name, module and status badges, field definitions table (FieldName, FieldLabel, DataType, ValueHelpType, IsRequired columns), and three action buttons in the header.}

\paragraph{Step 3: Edit Configuration Lines}
The system creates a new request with Status=DRAFT and navigates to the module-specific configuration app (e.g., FI Limit on port 8084). The page opens in edit mode with the current DEV data loaded in the table.

Use the toolbar buttons to modify configuration lines:
\begin{itemize}[leftmargin=2em]
  \item \textbf{Add Line:} Inserts a new blank row at the top of the table
  \item \textbf{Import from Excel:} Opens a file picker to upload an Excel file (.xls, .xlsx, .csv) with configuration data. The system maps column headers automatically using case-insensitive aliases.
  \item \textbf{Delete:} Marks selected rows for deletion (shown with red highlight)
\end{itemize}

Edit field values directly in the table cells. Value help dialogs are available for fields like ExpenseType, GlAccount, Currency, Plant, Warehouse, etc. Row highlights indicate the change type: green for new rows, blue for modified rows, red for deleted rows.

\imgplaceholder{INSERT: Screenshot --- Config app in edit mode. Shows DynamicPage with header cards (Request ID and Status, Configuration name, Title input and Reason textarea). Summary statistics strip (Total, Create, Update, Delete counts). Table with editable Input fields, value help buttons, and row highlights. Toolbar with Add Line, Import Excel, Delete, Save Draft, and Cancel buttons.}

\paragraph{Step 4: Save Draft}
Click \textbf{``Save Draft''} to persist your changes. The system performs the following:
\begin{enumerate}[leftmargin=2em]
  \item Validates all rows (required fields, business rules such as SendWh $\neq$ ReceiveWh)
  \item Fetches a CSRF token from the OData service
  \item For each changed row: creates a draft entity via POST, then activates it
  \item Updates the request header with the title and reason
  \item Reloads the table with the updated overlay
\end{enumerate}

If validation errors are found, the system highlights the problematic fields with red borders and displays an error message listing all issues. The save operation is blocked until all errors are resolved.

\paragraph{Step 5: Submit Request}
Click \textbf{``Submit Request''} to send the request for manager approval. The system:
\begin{enumerate}[leftmargin=2em]
  \item Validates that the request has at least one item
  \item Shows a confirmation dialog
  \item Calls the \texttt{submit} bound action on the backend
  \item Sends an email notification to relevant stakeholders
  \item Changes the status to SUBMITTED
  \item Navigates back to the request list
\end{enumerate}

\imgplaceholder{INSERT: Screenshot --- Submit confirmation dialog asking the user to confirm submission, followed by a success message: ``Request submitted for approval.''}

\paragraph{Common Mistakes and Tips}

Users should be aware of the following common mistakes when creating and submitting configuration requests:

\begin{itemize}[leftmargin=2em]
  \item \textbf{Forgetting to save before submitting:} If you add or modify configuration lines but do not click ``Save Draft'' before clicking ``Submit Request,'' the unsaved changes will not be included in the submission. Always save your draft first, verify the changes in the table, and then submit.
  \item \textbf{Duplicate business keys:} Each configuration module has a composite business key (e.g., ExpenseType + GlAccount for FI Limit). If you add two rows with the same key combination, the save operation will flag a duplicate error (BR-013). When importing from Excel, check for duplicate keys in the spreadsheet before importing.
  \item \textbf{Incorrect date formats in Excel:} For SD Price configurations, the ValidFrom and ValidTo columns in the Excel file should contain recognizable date values. The parser handles ISO 8601 (yyyy-MM-dd), standard locale formats, and Excel serial date numbers, but free-text date strings (e.g., ``next Monday'') will not be parsed correctly and will result in empty date fields after import.
  \item \textbf{Using Excel import for bulk changes:} When you need to update more than 5--10 configuration lines, using the ``Import from Excel'' feature is significantly faster than adding lines one by one. Prepare your changes in a spreadsheet with the correct column headers, import the file, review the imported rows in the table (they appear with green highlighting), and then save the draft. The import feature supports up to 500 rows without a confirmation prompt; larger imports trigger a confirmation dialog to prevent accidental mass data entry.
  \item \textbf{Checking the changelog before submitting:} Before submitting a request, navigate to the Catalog Landing page for the same configuration and use the ``Show Changelog'' button on existing requests to review the history of previous changes. This helps avoid submitting redundant changes or changes that conflict with recently approved requests.
\end{itemize}


\subsection{Workflow 2: Approve or Reject Request}

\textbf{Purpose:} Managers review submitted requests and either approve or reject them, with full visibility into the proposed configuration changes.

\paragraph{Step 1: Open Manager Dashboard}
Navigate to the Manager App. The Dashboard displays animated KPI tiles (Pending, Total, Approved, Rejected), a status donut chart, a module column chart, and a table of recent pending approvals.

\imgplaceholder{INSERT: Screenshot --- Manager Dashboard. Shows 4 KPI tiles with RadialMicroChart indicators, two VizFrame charts (Status Distribution donut and Requests by Module column chart), and a recent pending table with clickable rows.}

\paragraph{Step 2: Review Request}
Click on a pending request to open its Object Page. The page shows:
\begin{itemize}[leftmargin=2em]
  \item Request header with title, creator, and status badge
  \item Configuration context (module, environment, configuration name)
  \item Approval decision section (status, reason, approver/rejector information)
  \item Items table with a \textbf{``Compare''} button per row to view old vs.\ new values side-by-side
\end{itemize}

\imgplaceholder{INSERT: Screenshot --- Manager Object Page showing the request header with green Approve and red Reject buttons, configuration context, and items table with Compare buttons.}

\paragraph{Step 3: Approve or Reject}

\begin{itemize}[leftmargin=2em]
  \item \textbf{To Approve:} Click the green \textbf{``Approve''} button. The backend validates all items (required fields, no duplicates, source record exists), creates an approval snapshot for potential rollback, and changes the status to APPROVED. An email notification is sent to the IT Admin team.
  \item \textbf{To Reject:} Click the red \textbf{``Reject''} button. Enter a rejection reason in the prompt dialog (mandatory). The status changes to REJECTED (terminal state), and the request creator is notified via email.
\end{itemize}

\paragraph{Tips for Efficient Review}

Managers reviewing configuration change requests should follow these practices for efficient and accurate review:

\begin{itemize}[leftmargin=2em]
  \item \textbf{Always use the Compare button:} For each item in the request, click the ``Compare'' button to see a side-by-side view of old and new values. Changed fields are highlighted with red (old) and green (new) backgrounds, making it easy to identify exactly what is being modified. This is especially important for Update (ActionType=``U'') items where only some fields may have changed.
  \item \textbf{Check the request reason field:} The Key User is prompted to provide a reason when creating the request. Review this field on the Request Object Page header section to understand the business justification for the proposed changes before examining the individual items.
  \item \textbf{Provide specific rejection reasons:} When rejecting a request, provide a specific and actionable reason rather than a generic message. For example, instead of ``Changes not approved,'' write ``GL Account 600200 is not valid for the TRAVEL expense type --- please use GL Account 600100 instead.'' This helps the Key User address the issue quickly if they choose to resubmit the request.
  \item \textbf{Use the Dashboard for prioritization:} The Manager Dashboard displays the count of pending requests and a table of the most recent submissions sorted by creation date. Use this to prioritize which requests to review first, especially when multiple requests are pending approval simultaneously.
\end{itemize}


\subsection{Workflow 3: Promote and Rollback}

\textbf{Purpose:} IT Administrators manage the promotion of approved configurations through the DEV $\rightarrow$ QAS $\rightarrow$ PRD environment pipeline, with rollback capability at each stage.

\paragraph{Step 1: Open IT Admin Inbox}
Navigate to the IT Admin App. The FlexibleColumnLayout (My Inbox pattern) displays:
\begin{itemize}[leftmargin=2em]
  \item \textbf{Left column (33\%):} Request list with KPI badges showing counts per environment (DEV, QAS, PRD), module tabs (MM/MMSS/FI/SD), and filter controls (search, environment segment, date range)
  \item \textbf{Right column (67\%):} Audit trail by default, or request detail panel when a request is selected
\end{itemize}

\imgplaceholder{INSERT: Screenshot --- IT Admin FCL layout. Left column shows request list with module tabs, KPI badges (5 DEV, 3 QAS, 2 PRD), and filter strip. Right column shows audit trail entries (APPROVE, PROMOTE, ROLLBACK actions with user, date, and module badges).}

\paragraph{Step 2: Select Request and Review}
Click on an approved request in the list. The detail panel opens in the right column with three tabs:
\begin{itemize}[leftmargin=2em]
  \item \textbf{Preview:} Table of all configuration changes, color-coded by action type (green=added, blue=modified, red=deleted), with a ``Show Changes Only'' toggle and version numbers
  \item \textbf{Journey:} Visual representation of the DEV $\rightarrow$ QAS $\rightarrow$ PRD promotion path, showing status badges, user, and date for each environment
  \item \textbf{Info:} Request metadata (title, reason, creator, dates)
\end{itemize}

\imgplaceholder{INSERT: Screenshot --- Detail panel with Journey tab active. Shows 3 environment cards (DEV with Active status and checkmark, QAS empty, PRD empty) connected by lines. Promote button in footer toolbar, Rollback button on the active node.}

\paragraph{Step 3: Promote Configuration}
Click the \textbf{``Promote''} button in the footer toolbar. The system:
\begin{enumerate}[leftmargin=2em]
  \item Determines the next environment (DEV $\rightarrow$ QAS, or QAS $\rightarrow$ PRD)
  \item Shows a confirmation dialog: ``Promote to QAS?''
  \item Copies configuration data from the current environment to the next environment for each module table
  \item If the target is PRD: increments the global version number
  \item Creates an audit log entry (action=PROMOTE)
  \item Sends an email notification
\end{enumerate}

The system enforces strict promotion order: DEV $\rightarrow$ QAS $\rightarrow$ PRD. Skipping environments is not allowed (BR-005).

\paragraph{Step 4: Rollback (if needed)}
If a problem is discovered after promotion, click the rollback button on the Journey card for the affected environment. The system:
\begin{enumerate}[leftmargin=2em]
  \item Checks the rollback order: PRD must be rolled back before QAS, QAS before DEV (BR-006)
  \item Shows a confirmation panel with the target environment
  \item Restores the previous configuration state from the audit log snapshot
  \item Creates a new audit log entry (action=ROLLBACK) with reversed old/new data
  \item Changes the status to ROLLED\_BACK
\end{enumerate}

\imgplaceholder{INSERT: Screenshot --- Rollback confirmation panel with orange/red border. Shows the target environment, a warning message, and Confirm Rollback / Cancel buttons.}

\paragraph{Step 5: Compare Environments}
Click the \textbf{``Compare''} button in the header to open the Environment Comparison dialog. Select a module from the dropdown and click ``Compare'' to see a side-by-side view of configuration values across DEV, QAS, and PRD.

The dialog displays a table with columns: Record Key, Field, DEV value, QAS value, PRD value, and Sync status. Values are color-coded: green (identical across environments), orange (different from majority), red (missing in one or more environments). A ``Diffs Only'' toggle filters to show only records with differences.

\imgplaceholder{INSERT: Screenshot --- Environment Compare dialog. Shows module dropdown, Compare button, Diffs Only toggle, and legend (green=Identical, orange=Different, red=Missing). Table with Record Key, Field, DEV, QAS, PRD columns with color-coded ObjectStatus cells. Version row per record group.}

\paragraph{Tips for IT Administrators}

IT Administrators managing the environment promotion pipeline should be aware of the following best practices:

\begin{itemize}[leftmargin=2em]
  \item \textbf{Review the Preview tab before promoting:} Before clicking ``Promote,'' always switch to the Preview tab to review the complete list of configuration changes that will be applied to the next environment. Use the ``Show Changes Only'' toggle to focus on modified records and verify that the expected changes are present.
  \item \textbf{Use Environment Compare after promotion:} After promoting a configuration to QAS or PRD, open the Environment Compare dialog and select the relevant module. Verify that the promoted values appear correctly in the target environment column and that the sync status indicators show green (identical) for the promoted records. This provides immediate confirmation that the promotion completed successfully.
  \item \textbf{Understand rollback ordering:} The system enforces strict rollback ordering (PRD first, then QAS, then DEV). If you need to roll back a configuration that has been promoted to all three environments, you must roll back PRD first, then QAS, then DEV. Attempting to roll back QAS while PRD still contains the promoted data will result in an error message identifying the blocking environment.
  \item \textbf{Check the Audit Trail regularly:} The Audit Trail panel (displayed by default in the right column when no request is selected) shows all recent actions across the system. Use this to monitor promotion and rollback activity, especially when multiple IT Administrators are working on different requests simultaneously. Each audit entry includes the action type, module, environment, user, and timestamp, providing a comprehensive activity log.
  \item \textbf{Use the Journey tab for status overview:} The Journey tab provides a visual timeline of the DEV to QAS to PRD promotion path for each request. The active environment is highlighted with a colored background and checkmark icon, making it easy to see at a glance where each request currently stands in the promotion pipeline.
\end{itemize}


% ============================================================
%                      REFERENCES
% ============================================================
\clearpage
\chapter*{References}
\addcontentsline{toc}{chapter}{References}

The following references were consulted during the development of this project. Citations follow the Harvard referencing style.

\begin{enumerate}[leftmargin=2em]

  \item SAP SE (2024) \textit{SAP Fiori Design Guidelines}. Available at: \url{https://experience.sap.com/fiori-design/} (Accessed: March 2026).

  \item SAP SE (2024) \textit{SAPUI5 SDK Documentation}. Available at: \url{https://sapui5.hana.ondemand.com/} (Accessed: February 2026).

  \item SAP SE (2024) \textit{RESTful ABAP Programming Model (RAP) --- Developer Guide}. Available at: \url{https://help.sap.com/docs/btp/sap-abap-restful-application-programming-model/} (Accessed: January 2026).

  \item SAP SE (2024) \textit{OData Version 4.0 Protocol --- SAP Gateway Foundation}. Available at: \url{https://help.sap.com/docs/SAP_GATEWAY} (Accessed: February 2026).

  \item SAP SE (2024) \textit{CDS View Entities --- ABAP CDS Development Guide}. Available at: \url{https://help.sap.com/docs/abap-cloud/abap-rap/cds-view-entities} (Accessed: January 2026).

  \item SAP SE (2024) \textit{SAP Fiori Elements --- Feature Showcase}. Available at: \url{https://github.com/niclas-gsp/sap-fiori-elements-feature-showcase} (Accessed: March 2026).

  \item SAP Community (2024) \textit{abapGit Documentation}. Available at: \url{https://docs.abapgit.org/} (Accessed: January 2026).

  \item SheetJS LLC (2024) \textit{SheetJS Community Edition Documentation}. Available at: \url{https://docs.sheetjs.com/} (Accessed: February 2026).

  \item Render Inc. (2024) \textit{Render Documentation --- Web Services}. Available at: \url{https://render.com/docs/web-services} (Accessed: March 2026).

  \item Google LLC (2024) \textit{Nodemailer --- Send emails from Node.js}. Available at: \url{https://nodemailer.com/} (Accessed: February 2026).

  \item IEEE (2011) \textit{IEEE Std 830-1998 --- Recommended Practice for Software Requirements Specifications}. IEEE Computer Society.

  \item Sommerville, I. (2015) \textit{Software Engineering}. 10th edn. Harlow: Pearson Education.

  \item Pressman, R.S. and Maxim, B.R. (2020) \textit{Software Engineering: A Practitioner's Approach}. 9th edn. New York: McGraw-Hill Education.

  \item SAP SE (2024) \textit{Managed Business Objects with Draft --- RAP Guide}. Available at: \url{https://help.sap.com/docs/btp/sap-abap-restful-application-programming-model/draft} (Accessed: February 2026).

  \item OASIS (2014) \textit{OData Version 4.0 --- OASIS Standard}. Available at: \url{https://www.oasis-open.org/standards/} (Accessed: January 2026).

\end{enumerate}


% ============================================================
%                    END OF DOCUMENT
% ============================================================

\vfill
\begin{center}
\rule{0.5\textwidth}{0.4pt}\\[6pt]
\textit{--- End of Document ---}
\end{center}

\end{document}
